\documentclass[11pt,oneside]{article}
\usepackage[T1]{fontenc}
\usepackage[utf8]{inputenc}
\usepackage{lmodern,microtype}
\usepackage[a4paper,margin=27mm,headheight=14pt]{geometry}
\usepackage{amsmath,amssymb,amsthm,mathtools}
\usepackage{booktabs,longtable,array,enumitem}
\usepackage{xcolor,xurl,hyperref,fancyhdr,listings}
\definecolor{ink}{HTML}{17324D}
\definecolor{accent}{HTML}{246B75}
\hypersetup{colorlinks=true,linkcolor=ink,citecolor=accent,urlcolor=accent,
 pdftitle={Coherent finite methods for intuitionistic primitive recursive arithmetic},
 pdfauthor={Research report prepared for Vladimir Reshetnikov}}
\pagestyle{fancy}\fancyhf{}
\fancyhead[L]{\small Coherent finite methods in iPRA}
\fancyhead[R]{\small Research report}
\fancyfoot[C]{\thepage}
\setlength{\parindent}{0pt}\setlength{\parskip}{5pt}
\setlength{\emergencystretch}{3em}
\setlist{topsep=4pt,itemsep=3pt}
\newcolumntype{P}[1]{>{\raggedright\arraybackslash}p{#1}}
\newtheorem{theorem}{Theorem}[section]
\newtheorem{lemma}[theorem]{Lemma}
\newtheorem{proposition}[theorem]{Proposition}
\newtheorem{corollary}[theorem]{Corollary}
\theoremstyle{definition}
\newtheorem{definition}[theorem]{Definition}
\newtheorem{example}[theorem]{Example}
\theoremstyle{remark}\newtheorem{remark}[theorem]{Remark}
\newcommand{\ipra}{\ensuremath{\mathrm{iPRA}}}
\newcommand{\N}{\mathbb N}\newcommand{\Z}{\mathbb Z}
\newcommand{\Q}{\mathbb Q}\newcommand{\R}{\mathbb R}
\newcommand{\PR}{\mathrm{PR}}
\newcommand{\Valid}{\operatorname{Valid}}
\newcommand{\CertCheck}{\operatorname{Check}}
\newcommand{\Run}{\operatorname{Run}}
\newcommand{\wt}{\operatorname{wt}}
\newcommand{\Spec}{\operatorname{Spec}}
\newcommand{\rank}{\operatorname{rank}}
\newcommand{\id}{\operatorname{id}}
\DeclareMathOperator{\lc}{lc}
\lstset{basicstyle=\ttfamily\footnotesize,breaklines=true,
 columns=fullflexible,keepspaces=true,frame=single,rulecolor=\color{black!20},
 xleftmargin=4pt,xrightmargin=4pt}
\begin{document}
\begin{titlepage}
\vspace*{14mm}
{\color{accent}\large FOUNDATIONS AND WORKED APPLICATIONS\par}
\vspace{8mm}
{\color{ink}\Huge\bfseries Coherent finite methods\par}
\vspace{4mm}
{\LARGE for intuitionistic primitive\par recursive arithmetic\par}
\vspace{10mm}
{\large A reusable framework for refinement, lifting,\par and finite algebraic certificates\par}
\vspace{10mm}
Prepared for Vladimir Reshetnikov\par
15 November 2026\par
\vspace{10mm}
\begin{minipage}{0.94\textwidth}
\textbf{Main results.} A finite filtered-algebra lifting theorem produces
unique compatible solutions in polynomial jets and residue rings by an
explicit primitive recurrence. A weighted determinant norm gives a sharp
finite point-count theorem for a class of plane curves, replacing a
divisor-theoretic input by exact polynomial arithmetic. A common proof
discipline connects these constructions to quantifier-free induction and
to ProveIt's actual object-language proof relation.

\textbf{Status.} The report gives complete human-readable arguments for its
stated finite theorems and explicit induction audits. Selected foundational
Lean proofs and concrete definitions are compiled. The full object-iPRA
derivations of the lifting and norm theorems remain formalization work;
their precise targets and completion steps are supplied.
\end{minipage}
\vfill
{\small Source snapshot: ProveIt\par
\texttt{052fbaebfc955294b4c02f6532ab7e7a3f760f7a}\par
Existing framework packets are treated as prior work, with additions
identified explicitly.}
\end{titlepage}
\begin{abstract}
Embedding a proof into first-order intuitionistic primitive recursive
arithmetic requires more than an executable algorithm. It requires finite
representations, primitive recursive constructions, and a derivation whose
induction invariants are decidable bounded predicates. We organize these
requirements around coherent finite presentations: every precision is a
finite arithmetic problem, restriction maps preserve its meaning, and
finite uniqueness supplies compatibility without countable choice.

The first main application is a division-free correction iteration in a
nilpotent filtered algebra. A supplied inverse Jacobian modulo the first
filtration level gives exact solutions after a stated number of steps.
Ordinary formal power series, simultaneous implicit equations, and roots
modulo powers of any integer base become instances of the same theorem.
A related Hurwitz-series presentation treats polynomial differential
equations without dividing by factorials. The second application gives
exact determinant degrees and point counts on curves with leading terms
\(Y^m-cX^d\), where \(\gcd(m,d)=1\). Fiber multiplicities are handled by
Vandermonde row certificates, so distinct points with a common abscissa
are counted correctly, including singular fibers.

These results address concrete representation and geometric counting
obligations in the inspected repository. They do not supply a general
height bound, a rational-point enumeration cutoff, analytic convergence
for an arbitrary formal solution, or the major analytic estimates absent
from the prime-distribution reports. The accompanying source audit,
exact-arithmetic programs, Lean files, and proof-status ledger distinguish
each of these boundaries.
\end{abstract}
\tableofcontents
\newpage
\section{Source audit, aim, and contribution}
\label{sec:audit}

\subsection{The inspected snapshot}

The repository was read through its GitHub connector at commit
\begin{center}\small\ttfamily
052fbaebfc955294b4c02f6532ab7e7a3f760f7a.
\end{center}
The proof-packet index lists 54 packets grouped by mathematical area.
It explicitly says that placing them in the repository did not verify,
rerun or compile their contents \cite{repoIndex}. We therefore distinguish
an author's stated result, a mathematical argument read in a source file,
a compiled host Lean theorem, and an actual object-iPRA derivation.

The review covered the index, the three previous framework packets, the
ConciseIPRA interfaces needed for the new Lean files, and selected
application reports. It is a targeted source audit, not a verification of
all 54 packets. The machine-readable manifest records inspected paths,
the pinned revision, available Git blob identifiers, and SHA-256 hashes
of materialized source text. Exact audit notes are included separately.

Three framework packets already exist at this revision:
\begin{enumerate}
\item \texttt{ipra\_framework\_2026-11-15}: an earlier compact framework.
\item \texttt{ipra\_embedding\_framework}: finite-demand transfer,
approximation and isolation, analytic regular branches, stable windows,
the odd-degree hyperelliptic norm count, and finite-algebra certificates.
\item \texttt{ipra\_framework}: finite contracts, an expanded analytic
chart constructor, finite-window arguments, and integral power-gap
certificates with explicit producers.
\end{enumerate}
Their reusable ideas are retained and cited as prior work
\cite{repoSmallFramework,repoEmbedding,repoFramework}. In particular,
bounded Booleanization, certificate composition, scalar Hensel steps,
ordinary bounded descent, and metastability are not presented as newly
discovered results of this report.

\subsection{The obstacles that determine the design}

\begin{longtable}{P{.20\linewidth}P{.34\linewidth}P{.37\linewidth}}
\toprule
Source & Specific difficulty or already available result & Consequence here\\
\midrule\endhead
ConciseIPRA & The equivalence transports an existing arithmetic derivation;
it does not turn arbitrary ambient Lean truth into one. & Every formal
endpoint is stated using \texttt{Prf PRAAx}.\\[4pt]
Existing frameworks & Finite observation contracts and bounded loops are
already developed. & Add exact restriction and coherence interfaces with
new worked algebraic constructions.\\[4pt]
Puiseux & General exact branch selection is impossible uniformly from
arbitrary Cauchy names; exact algebraic and regular-chart results exist.
& Solve simultaneous regular formal equations from a supplied invertible
Jacobian, with exact coherent jets. Retain the singular obstruction.\\[4pt]
Lagrange inversion; Bell numbers & Many identities concern finite
coefficient lists even when written using infinite generating functions.
& Give a common formal lifting and integral Hurwitz-series interface;
use their known identities as illustrations.\\[4pt]
Siegel general route & The proper-subspace fiber bound is described via
functions, divisors, and a zero/pole degree theorem. & Construct functions,
an exact dimension count and a sharp zero-count certificate for weighted
plane curves using finite polynomials.\\[4pt]
Earlier embedding framework & The fiber repair is already proved for
\(Y^2=f(X)\) with odd \(\deg f\). & Extend to \(Y^m-cX^d\) with arbitrary
lower-weight mixed terms and \(\gcd(m,d)=1\).\\[4pt]
Mordell--Weil & Contracting descent with explicit fuel already exists; actual
quotient representatives remain an input. & Reuse bounded descent only
as infrastructure. Do not claim it constructs those representatives.\\[4pt]
Faltings; Siegel & A general cardinal bound, height bound and complete list
are different targets. & Supply a counting component and state the
unresolved cover and height obligations separately.\\[4pt]
Diophantine geometry & Modular obstructions and bounded searches are finite;
passing local tests does not establish global rational points. & Add
compatible lifting through modulus powers; assert no local-to-global
converse.\\[4pt]
Prime distribution & Field certificates and major analytic estimates are
separate inputs. & The finite algebra modules can support certificates;
they do not prove the missing analytic estimates.\\
\bottomrule
\end{longtable}

The application rows refer to the inspected Puiseux, Lagrange-inversion,
Bell-number, Mordell--Weil, Faltings, Diophantine-geometry and prime-distribution
packets \cite{repoPuiseux,repoLagrange,repoBell,repoMW,repoFaltings,repoGeometry,repoPrimes}.

The most concrete source obligation addressed below is the paragraph
``The geometric input to the fibre bound'' in the Siegel report's
\texttt{paper/05\_general\_route.tex} \cite{repoSiegel}. Its original
description requests finite functions, independence data and a zero-count
argument. The earlier hyperelliptic repair settles a restricted instance.
Our weighted norm theorem settles a broader, explicitly stated family
without assuming a developed divisor theory. It does not settle the
separate finite-tuple approximation theorem required by that route.

\subsection{What this report adds}

The central framework addition is \emph{coherent finite presentation}.
Instead of immediately constructing an infinite object, construct its
requested finite observation, prove its finite specification, and prove
that restriction preserves that specification. A finite uniqueness theorem
then supplies coherence. This allows exact algebra, modular arithmetic and
formal differential algebra to use the same proof architecture as the
repository's metric approximation modules, while preserving their different
meanings.

The first main application is a filtered-algebra lifting theorem. For a
polynomial system, a root and an inverse Jacobian modulo the first
filtration level yield a fixed correction recurrence. Each step raises the
order of the residual by one; an explicitly nilpotent truncation therefore
needs only a prescribed finite number of steps. This gives formal implicit
functions in several unknowns and roots modulo powers of composite as well
as prime integer bases. A companion Hurwitz construction removes factorial
division from polynomial differential equations.

The second main application is an exact norm theorem for weighted plane
curves. It proves the determinant's degree and leading coefficient, then
shows that a fiber containing \(r\) distinct common zeros contributes
\(r\) factors to that determinant. The result gives a sharp point count,
an elementary dimension formula for an explicit space of functions, and a
proper-subspace fiber bound. An adjugate argument also gives an exact
function-field presentation for this curve class.

These use established algebraic ideas. Formal implicit-function theorems
and constructive Hensel methods have substantial prior literature
\cite{Hu,AlonsoLombardiNeuwirth}; norms as multiplication determinants are
standard \cite{StacksNorm}. The contribution claimed here is the explicit
generalization relative to the inspected packets, the integrated finite
proof architecture and its arithmetic audit. No priority claim over the
mathematical literature is made.

\subsection{Reading and proof-status conventions}

Sections~\ref{sec:foundations} and \ref{sec:refinement} give the common
rules. The next two chapters contain the main mathematical applications.
Section~\ref{sec:formalization} maps them to concrete Lean definitions,
compiled lemmas and remaining object derivations. The final chapter lists
exact accomplishments and boundaries.

The phrase ``has an iPRA proof'' in a mathematical theorem means that the
following human argument is expressed by the stated PR constructions and
bounded induction invariants. It does not mean that a corresponding Lean
proof object has already been checked. A compiled host theorem verifies
its own Lean proposition. A compiled theorem whose conclusion is
\texttt{Prf PRAAx ...} verifies an actual object derivation or a transformation
of supplied object derivations. These statuses are recorded individually;
compiling an unproved proposition definition is never counted as proving it.

\section{The arithmetic target and its proof discipline}
\label{sec:foundations}

\subsection{The theory used throughout}

The object theory is first-order intuitionistic arithmetic with natural
numbers as its only sort. Its language has \(0,S\) and a symbol for every
fixed primitive recursive construction. The symbols are generated from
zero, successor and projections by composition and primitive recursion.
For a recursion symbol the defining equations are
\[
 r(0,\boldsymbol x)=g(\boldsymbol x),\qquad
 r(n+1,\boldsymbol x)=h(n,r(n,\boldsymbol x),\boldsymbol x).
\]
We include the successor axioms \(Sx\ne0\) and
\(Sx=Sy\to x=y\), equality reflexivity and substitution, and ordinary
intuitionistic first-order natural deduction. In particular implication,
conjunction, disjunction, falsity and both quantifiers have their usual
introduction and elimination rules, with the standard freshness conditions
on quantifier introduction and existential elimination. No excluded-middle
schema, choice principle or Markov principle is added.

Induction is restricted to quantifier-free formulas, with parameters:
\[
 A(0,\boldsymbol x),\quad
 \forall n\bigl(A(n,\boldsymbol x)\to A(n+1,\boldsymbol x)\bigr)
 \quad\Longrightarrow\quad \forall n\,A(n,\boldsymbol x).
\]
This is the ordinary \ipra\ presentation formalized by
\texttt{PRAAx}, \texttt{PTm}, and \texttt{PR} in ConciseIPRA
\cite{repoCore}. The native concise calculus is connected to it by an
existing syntactic bi-interpretation. We work first in this explicit
arithmetic presentation, and transport a derivation only after it exists.

The induction restriction does not prohibit formulas with quantifiers from
appearing as assumptions or conclusions. For example, a hypothesis
\(\forall n\,C(n)\) may be instantiated at each of finitely many indices in
a calculation. What must be audited is the formula on which induction is
performed.

\subsection{Finite objects, total operations, and decidability}

An integer is a signed natural-number code; a rational is a pair
\((z,d)\) with \(d>0\). We do not need quotient objects: equality is
\((z,d)\simeq(w,e)\) precisely when \(ze=wd\). Addition, multiplication,
comparison and inversion with a nonzero-numerator guard are explicit
integer calculations. A gcd normalization is optional. All operations
receive harmless default values on malformed input; the theorem carries
the validity guards.

Finite tuples, coefficient arrays, matrices and execution histories are
also natural-number codes. One possible list code is empty list \(0\)
and
\[
 \operatorname{cons}(a,t)=1+2^a(2t+1).
\]
Recovering its head and tail uses bounded tests on a code; the tail is
smaller than a nonempty valid code. Repeated decoding can therefore be
implemented with a loop bounded by that code. This very inefficient coding
is nevertheless primitive recursive. A practical implementation may use
a more efficient pairing scheme; no argument depends on efficiency.

A polynomial is a finite coefficient list or exponent-coefficient table.
Multiplication is a bounded convolution. A determinant is the signed sum
over finitely many permutations, with at most \(n!\) summands. A rational
matrix can be row-reduced by scanning bounded lists of columns and rows;
each pivot is selected by a decidable nonzero test. These descriptions
establish that the operations are primitive recursive. Their algebraic
correctness still needs proofs, given by the finite invariants described
below. A computation model is not a correctness argument.

\begin{lemma}[Booleanization, as a schema]
\label{lem:booleanization}
For every fixed bounded arithmetic formula \(B(\boldsymbol x)\), built
from fixed PR terms and bounded number quantifiers, there is a fixed PR
function \(b_B\) with values in \(\{0,1\}\), and an \ipra\ derivation
\[
 B(\boldsymbol x)\ \longleftrightarrow\ b_B(\boldsymbol x)=1.
\]
Consequently induction on bounded formulas is a derived use of
quantifier-free induction.
\end{lemma}
\begin{proof}
Equality and order have their ordinary PR Boolean tests. The successor
axioms and quantifier-free induction prove their correctness. Boolean
complement, product and maximum implement negation, conjunction and
disjunction. The Boolean implication test returns one exactly when the
antecedent test is zero or the consequent test is one. Its equivalence to
intuitionistic implication uses the already proved decidability of the
antecedent; it does not decide arbitrary propositions.

For a bounded universal quantifier use a prefix product
\(p(0)=1\), \(p(t+1)=p(t)c(t)\), where \(c(t)\in\{0,1\}\).
Quantifier-free induction proves \(p(t)\in\{0,1\}\) and
\(p(t+1)\le p(t)\). Fix an arbitrary index \(i\). Induction on the
quantifier-free implication
\(i<t\land p(t)=1\to c(i)=1\) proves the forward direction of
\(p(b)=1\leftrightarrow\forall i<b\,c(i)=1\).
For the reverse direction, keep the bounded universal assertion as a
fixed assumption and induct on \(t\le b\to p(t)=1\). In the step
\(t<b\), instantiate that assumption at \(t\).

For bounded existence scan all positions below the bound, retaining the
first success and using the current scan length as a failure sentinel
(so the final sentinel is the final bound). Induct
on the conjunction of: the stored value is at most the current length;
if it is smaller it is a successful witness; and the prefix product of
the failure tests before it is one. These are quantifier-free tests.
A supplied successful position precludes the sentinel, and the stored
position then witnesses the existential assertion. This proves both
directions without double-negation elimination.

Structural induction in the metatheory on the fixed formula supplies
the test and its derivation. Replacing a bounded induction formula by
the equation \(b_B(n,\boldsymbol x)=1\) completes the last assertion.
\end{proof}

This proof also gives bounded minimization and finite choice. In particular,
from a proof of
\(\forall i<n\,\exists y<B(x,i)\,R(x,i,y)\), with fixed PR bounds and
tests, a PR scan produces the finite list of least witnesses. There is no
corresponding inference from unbounded existence without an additional
bound or a separately proved PR constructor.

\subsection{A precise finite-loop rule}

Let the fixed PR run be
\[
 S(x,0)=\operatorname{init}(x),\qquad
 S(x,t+1)=\operatorname{step}(x,t,S(x,t)).
\]
Let \(P(x)\), \(I(x,N,t,s)\), and \(Q(x,N,s)\) be bounded predicates.
Suppose there are already \ipra\ proofs of
\begin{align*}
 P(x)&\to I(x,N,0,\operatorname{init}(x)),\\
 t<N\land I(x,N,t,s)&\to
 I(x,N,t+1,\operatorname{step}(x,t,s)),\\
 I(x,N,N,s)&\to Q(x,N,s).
\end{align*}
Then \ipra\ proves \(P(x)\to Q(x,N,S(x,N))\). Indeed, under the
fixed assumption \(P(x)\), induct on the Booleanization of
\[
              t\le N\ \longrightarrow\ I(x,N,t,S(x,t)).
\]
For the successor case with \(t+1\le N\), apply the step proof; outside
the guard the implication is immediate. Instantiating at \(N\) gives
the conclusion. This is the precise meaning of a finite induction audit
in later chapters.

The rule also accommodates a finite log of hypotheses queried from a
represented input. Include a Boolean conjunction saying that every
previously queried fact holds. A global regularity assumption can supply
the next fact by universal elimination, or the entire finite log can be
part of the theorem's input. The induction invariant remains a bounded
predicate about the coded state.

\begin{theorem}[Bounded descent with a reconstruction invariant]
\label{thm:rank}
Fix PR functions \(T,s_0,r\), a Boolean terminal test \(D\), and a
bounded validity predicate \(V\). Suppose \ipra\ proves that \(T\)
preserves \(V\), fixes terminal valid states, and satisfies
\[
 V(s)\land D(s)=0\quad\Longrightarrow\quad r(T(s))<r(s).
\]
If \(V(s_0)\), then the iterate \(T^{r(s_0)}(s_0)\) is terminal.
Every bounded invariant proved to be preserved by \(T\) holds there.
\end{theorem}
\begin{proof}
The bounded-loop rule gives validity of the iterates. Induct on
\[
 D(T^t(s_0))=1\quad\lor\quad r(T^t(s_0))+t\le r(s_0).
\]
If the terminal alternative holds, absorption preserves it. Otherwise
strict decrease of the natural-number rank preserves the second
alternative at \(t+1\). At \(t=r(s_0)\), a nonterminal state would
have rank zero. Its stipulated next step would have negative rank, a
contradiction. Hence the terminal test is one. A simultaneous bounded
invariant is proved by the same loop rule.
\end{proof}

This is background infrastructure, not a new Mordell--Weil argument.
The inspected Mordell--Weil packet already proves a more efficient
contracting-height descent once actual quotient representatives are
supplied. The theorem is useful here because a progress measure, a
coefficient precision and a finite row-reduction phase are all explicit
natural-number budgets. It also has a compiled Lean proof.

\subsection{What is uniform, and what is only a schema}

The degrees, matrix dimensions, integer moduli, coefficient lists and
precision bounds in the applications may be variable natural-number
inputs. Their algorithms use a fixed finite vocabulary of integer and
rational operations, together with explicitly bounded loops. They are
therefore genuinely uniform PR constructions.

An abstract coefficient ring is different. A theorem over every
\emph{PR-presented ring} is a schema: fix PR operations and supply or prove
their ring laws; then compile that instance. This does not introduce a
variable ranging over all total PR programs with a universal evaluator.
If such a PR evaluator \(U(e,x)\) existed, listing all unary PR functions,
then \(d(n)=U(n,n)+1\) would itself be PR and have an index \(e\).
Substitution of \(e\) gives \(U(e,e)=U(e,e)+1\), impossible. The
diagonal argument also explains why a generic evaluator hidden in an
apparently convenient Lean record would be a serious error.

Finally, a PR proof checker for arbitrary arithmetic proofs is not an
internal reflection principle. We never infer a formula from the single
assumption that an arbitrary coded proof passes a checker. We construct
derivations of the specific fixed formulas at issue. The object
composition lemmas in the archive transform actual \texttt{Prf PRAAx}
objects, with every input proof retained as a hypothesis.

\section{Coherent finite presentations}
\label{sec:refinement}

The existing packets use finite observations and explicit error contracts.
For algebraic constructions there is a useful further organizing principle:
successive finite observations can be \emph{exact restrictions} of one
another. This avoids treating all infinite objects as analytic Cauchy
sequences and exposes the algebra actually needed by the proof.

\subsection{A presentation at each requested precision}

For fixed numerical parameters \(x\), let \(A_n\) be a presentation of
finite data at precision \(n\). In the arithmetic theory this means a
bounded validity test \(V(x,n,z)\), a bounded equivalence test
\(z\simeq_{x,n}z'\), and PR restriction operations
\(\pi_{m,n}\) for \(n\le m\). Their laws, on valid inputs, are
\[
 \pi_{n,n}(z)\simeq_{x,n}z,\qquad
 \pi_{k,n}(\pi_{m,k}(z))\simeq_{x,n}\pi_{m,n}(z)
 \quad(n\le k\le m).
\]
All operations respect the equivalence tests and preserve validity.
For polynomials modulo \(T^{n+1}\), restriction deletes high
coefficients. For residues modulo \(q^n\), it reduces to the smaller
modulus. For finite matrix identities, precision may instead mean the
finite list of entries or identities requested.

A represented object is a specified PR producer \(a(x,n)\), together
with proofs of its validity and coherence
\[
 \pi_{m,n}(a(x,m))\simeq_{x,n}a(x,n)\qquad(n\le m).
\]
The first-order theorem quantifies over the numerical parameters and the
requested precisions. It does not assert the existence of a function
object inside \ipra. One can interpret the family externally as an
element of an inverse limit, but that interpretation is explanatory, not
an additional axiom used to establish the finite claims.

\begin{theorem}[Finite uniqueness supplies coherence]
\label{thm:coherence}
Let \(C(x,n,z)\) be a bounded specification that includes validity at
precision \(n\). Suppose a fixed PR producer
\(a(x,n)\) satisfies it, restrictions preserve it, and it determines a
unique equivalence class at every precision:
\begin{align*}
 &C(x,n,a(x,n)),\\
 &n\le m\land C(x,m,z)\to C(x,n,\pi_{m,n}(z)),\\
 &C(x,n,z)\land C(x,n,z')\to z\simeq_{x,n}z'.
\end{align*}
If these implications have \ipra\ derivations, so does the coherence
formula for \(a\).
\end{theorem}
\begin{proof}
Fix \(x,n,m\) with \(n\le m\). Correctness at \(m\) and restriction
give \(C(x,n,\pi_{m,n}(a(x,m)))\). Correctness at \(n\) gives
\(C(x,n,a(x,n))\). Apply finite uniqueness. No induction and no choice
principle is required in this final step.
\end{proof}

The statement deliberately includes existence \emph{with a fixed PR
producer}. A bare theorem \(\forall n\,\exists z\,C(n,z)\), even with
uniqueness, does not by itself provide a PR function. The lifting theorem
in the next chapter constructs the producer before coherence is invoked.
Moreover the specification includes the selected initial residue or
constant term. Without that seed there may be several roots and the
uniqueness hypothesis would be false.

\subsection{Operations carry an explicit precision demand}

An operation producing an output observation at level \(n\) must state
which input levels it reads. For product of ordinary jets the demand is
\(n\) for each input. For a derivative it is \(n+1\). For a coefficient
in an implicit solution the demand is the polynomial data and the seed
plus its finite inverse certificate. Once these demands are stated,
composition is ordinary composition of their PR bounds and programs.

Suppose a first operation has witness \(f(x)\) and proof
\(P(x)\to Q(x,f(x))\), and a second operation has witness \(g(x,y)\)
and proof \(Q(x,y)\to R(x,y,g(x,y))\). Substitution and implication
elimination give
\[
 P(x)\to R(x,f(x),g(x,f(x))).
\]
This is a syntactic proof transformation, not an appeal to semantic
soundness. The module \texttt{ObjectComposition.lean} implements it for
the repository's actual \texttt{Prf PRAAx} relation, including capture-free
substitution and specialization of parameters.

Coherent operations can similarly be composed. Restrict the outer result,
use its restriction law, apply the restriction laws of its arguments,
and then use congruence. Every equality is a finite equality at a stated
level. No equality test for the entire infinite object is needed.

\subsection{Three kinds of precision, with different meanings}

\begin{center}\small
\begin{tabular}{P{.20\linewidth}P{.30\linewidth}P{.40\linewidth}}
\toprule
Presentation & Finite observation & What passing every level means externally\\
\midrule
Ordinary jets & Coefficients through degree \(n\), exact modulo
\(T^{n+1}\) & A formal power series identity; it does not supply a radius
of analytic convergence.\\[3pt]
Residue towers & A root modulo \(q^n\), with a selected residue modulo
\(q\) & A compatible \(q\)-adic solution; it need not be a rational or
integer solution.\\[3pt]
Metric names & Rational observations with explicit error bounds & A real
or complex object in a stated represented family; exact zero tests do not
follow from arbitrary finite approximations.\\
\bottomrule
\end{tabular}
\end{center}

These are different embeddings of mathematics. Their finite arithmetic
theorems are useful in their own right. A conversion between them needs a
separate proof: convergence estimates convert suitable jets to analytic
names, and a global arithmetic argument may sometimes turn local residue
information into a rational solution. Neither conversion is automatic.

\subsection{Certificates: validity, production, and scope}

For every certificate language record three separate statements:
\[
 \CertCheck(x,c)=1\to \Spec(x,c),\qquad
 \CertCheck(x,\operatorname{produce}(x))=1,
\]
and the exact input class on which the second statement is proved.
The specification itself must be an explicit arithmetic formula, not an
unexplained proof field whose contents already are the desired theorem.

In the lifting application, the initial certificate consists of a root
modulo the first filtration level and a matrix inverse there. Both are
finite checks. The algorithm then produces all requested refinements.
For rational jets a nonzero rational determinant lets finite elimination
construct the inverse. For residues, the inverse may be supplied or
found by a bounded search modulo \(q\), provided invertibility is
established. In the norm application the certificate is constructed
directly: form the multiplication matrix and its determinant. Its exact
degree is proved from the weighted leading terms. No external oracle
supplies the norm.

The later covering corollary has a different status. Given a list of
covering subspaces, it proves a point bound. Producing such a list for
the integral or rational points on an arbitrary curve remains a separate
mathematical obligation. The type of the output is kept visible precisely
so that a counting theorem cannot be mistaken for a height bound.

\section{Finite filtered algebras: a common lifting theorem}
\label{sec:filtered-lifting}

The existing Lagrange-inversion packet constructs coefficient prefixes by
fixed-point iteration, and the Puiseux packet reduces branches to simple-root
equations before applying a coefficient recurrence.  The Bell packet uses
finite exponential-series identities.  These are substantial existing
arguments; the purpose of this section is to expose a common reusable
arithmetic theorem, including simultaneous polynomial systems and modular
lifting \cite{repoLagrange,repoPuiseux,repoBell}. The classical Hensel principle is not new. The contribution here
is its organization as one division-free finite program, one bounded
invariant, and one coherence argument suitable for iPRA.

\subsection{The exact arithmetic interface}

Fix a primitive-recursively presented commutative ring $B$.  Its element
codes, validity predicate, equality test, and ring operations are fixed
primitive recursive data, and the ring and congruence laws have supplied
iPRA derivations.  The displayed equalities below mean equality in this
presentation; a canonical representation is convenient but unnecessary.
Let $N\ge1$ be a numerical precision parameter.  Assume that
$F^kB$ is a uniformly primitive recursive membership test, with proofs of
the following laws:
\begin{enumerate}
\item $F^0B=B$, $F^{k+1}B\subseteq F^kB$, and $F^NB=\{0\}$;
\item every $F^kB$ is an additive subgroup;
\item $uv\in F^{i+j}B$ whenever $u\in F^iB$ and $v\in F^jB$.
\end{enumerate}
One can set $F^kB=\{0\}$ for $k\ge N$.  Because $F^0B=B$, the last
condition also gives closure under multiplication by an arbitrary ring
element.  Matrices and vectors belong to a filtration level when each of
their finitely many entries does.

These are \emph{supplied numerical contracts}.  This definition does not
assert that membership in powers of an arbitrary ideal is decidable.
The ring and filtration may depend on numerical parameters, including
$N$, provided all operations and membership tests are jointly primitive
recursive and the same displayed laws are proved uniformly.  The concrete
jet and residue-ring examples below have this property.

A polynomial in $m$ variables is a finite list of terms
$(c,e_1,\ldots,e_m)$, interpreted as
$cX_1^{e_1}\cdots X_m^{e_m}$.  Missing coefficients are zero; duplicate
monomials are allowed.  Evaluation, formal partial derivatives, vector
evaluation, and finite matrix multiplication are primitive recursive
bounded folds.  No general-purpose primitive recursive program evaluator
is used.  The polynomials themselves, their dimensions, and their degrees
may be variable numerical inputs.

\begin{lemma}[Finite polynomial difference and Taylor identities]
\label{lem:filtered-taylor}
Let $P\in B[X_1,\ldots,X_m]^m$, and let $J(x)$ denote its Jacobian.
If $h\in(F^kB)^m$, then
\[
 P(x+h)-P(x)-J(x)h\in(F^{2k}B)^m.                 \tag{T}
\]
If $x-a\in(F^1B)^m$, then every entry of $J(x)-J(a)$ belongs to $F^1B$.
Both claims have uniform iPRA derivations from the presentation contracts.
\end{lemma}

\begin{proof}
Expand one monomial
$c\prod_j(x_j+h_j)^{e_j}$ by the finite binomial theorem.  The term using
no $h_j$ is its value at $x$.  The terms using exactly one $h_j$ have sum
equal to its formal derivative applied to $h$.  Every remaining term has
at least two factors among the $h_j$, counted with multiplicity, and
therefore belongs to $F^{2k}B$.  Multiplication by the other factors
preserves membership, and summing finitely many such terms does so too.
Adding the finitely many monomials proves (T).

For the second assertion apply the elementary finite identity
$u^r-v^r=(u-v)\sum_{j<r}u^jv^{r-1-j}$ to one variable at a time in
each monomial of a partial derivative.  The resulting difference is a
finite sum of multiples of the coordinates $x_j-a_j$, hence belongs to
$F^1B$.  The proof uses no division and is valid in every characteristic.

The binomial and telescoping identities are proved by induction on natural
exponents.  Their induction formulas are equalities of coded ring elements.
The further invariants assert that all terms in a finite accumulated list
have the stated membership; they are bounded formulas and have Boolean
quantifier-free forms.  Thus the finite expansion argument gives actual
iPRA proof instructions uniformly in the polynomial code, rather than
merely an appeal to a general algebra theorem in the ambient metatheory.
\end{proof}

\subsection{A frozen inverse is enough}

\begin{theorem}[Finite filtered lifting]
\label{thm:filtered-lifting}
Let $a\in B^m$ and $U\in B^{m\times m}$ satisfy
\[
 P(a)\in(F^1B)^m,\qquad
 I-J(a)U\in(F^1B)^{m\times m},\qquad
 I-UJ(a)\in(F^1B)^{m\times m}.                 \tag{H}
\]
Define the primitive recursive sequence
\[
 b_0=a,\qquad b_{r+1}=b_r-U P(b_r).            \tag{L}
\]
Then $b_{N-1}$ is a root of $P$, and is the unique root $x$ satisfying
$x-a\in(F^1B)^m$.  No derivative inverse is recomputed during the
iteration, and no division occurs after the data $a,U$ are supplied.

More precisely, for every $r<N$,
\[
 b_r-a\in(F^1B)^m,\qquad
 P(b_r)\in(F^{r+1}B)^m,\qquad
 b_{r+1}-b_r\in(F^{r+1}B)^m.                 \tag{I}
\]
The first two assertions suffice at $r=N-1$; the last assertion is also
meaningful there if one more iterate is formed.
\end{theorem}

\begin{proof}[Existence and the precision invariant]
Put $x=b_r$, $v=P(x)$, and $h=-Uv$.  Suppose $x-a\in F^1$ and
$v\in F^k$, where $k\ge1$.  Matrix multiplication preserves membership
in $F^k$, so $h\in F^k$.  In particular $x+h-a\in F^1$.
Lemma~\ref{lem:filtered-taylor} gives
\[
 P(x+h)=(I-J(a)U)v-(J(x)-J(a))Uv+R,
 \qquad R\in(F^{2k}B)^m.                    \tag{1}
\]
The first term belongs to $F^{k+1}$ by (H).  The second belongs to
$F^{k+1}$ because $J(x)-J(a)\in F^1$.  The remainder does as well:
$2k\ge k+1$ and the filtration is decreasing.  Thus one iteration
raises the certified residual precision by at least one level.

Initially the first two assertions of (I) hold by (H).  The preceding
argument proves the induction step, and the defining equation (L) proves
the increment assertion.  At $r=N-1$ the residual belongs to $F^NB=\{0\}$.
If $N=1$, no update is performed and the hypothesis already says $P(a)=0$.
The proof uses the right inverse congruence $J(a)U\equiv I$; the left
inverse congruence will be used for uniqueness.
\end{proof}

\begin{proof}[Uniqueness]
Let $x,y$ be roots with $x-a,y-a\in F^1$.  We show inductively that
$d=x-y\in F^k$ for $1\le k\le N$.  The case $k=1$ is immediate.
If $d\in F^k$, the Taylor identity at $y$ yields
\[
 0=P(x)-P(y)=J(y)d+R,\qquad R\in F^{2k}.
\]
Multiplying by $U$ and rearranging gives
\[
 d=(I-UJ(a))d-U(J(y)-J(a))d-UR.              \tag{2}
\]
Each term on the right lies in $F^{k+1}$, for exactly the three reasons
used in (1).  This proves the step.  At level $N$, $d=0$, hence $x=y$.
The induction is finite and contains no least-counterexample principle,
choice, field cancellation, or existence assertion about a completion.
\end{proof}

The linear precision gain is a deliberate safe bound.  It need not be
quadratic when $U$ is only an inverse modulo $F^1$.  For example, in
$B=\mathbb Z/q^N\mathbb Z$, take $P(X)=X-q$, $a=0$, and $U=1+q$.
The residual satisfies $P(b_r)=(-q)^{r+1}$.  For $r+1<N$ it is generally
nonzero modulo $q^N$, so the bound of $N-1$ updates is attained.

\paragraph{Why this is an iPRA theorem scheme.}
Fix the final precision $N$ before induction.  At stage $r$, use
\[
 r<N\ \longrightarrow\
 \bigl(\operatorname{ValidVec}(b_r)\wedge
       b_r-a\in F^1\wedge P(b_r)\in F^{r+1}\bigr)
                                                               \tag{3}
\]
as the readable invariant.  Its vector and matrix conditions range only
over their coded finite lengths.  Booleanization converts (3) to a
quantifier-free equation.  The program symbols in (L) are obtained by one
primitive recursion; matrix and polynomial operations are previously fixed
PR definitions.  The existence proof consequently compiles by
\texttt{PRA.qfInd}.  For uniqueness, fix $x,y,N$ as parameters and use the
Booleanization of $k\le N\to x-y\in F^k$, treating the step from
$k=0$ separately by the known congruence modulo $F^1$.
The contracts for the ring and
filtration remain explicit hypotheses or previously supplied iPRA lemmas.
The argument does not deduce object-theory provability merely from a
terminating Lean implementation.

\subsection{Coherence without selecting an infinite object}

\begin{corollary}[Transport of the unique lift]
\label{cor:filtered-coherence}
Suppose $\rho:B\to C$ is a fixed PR map between two PR-presented
filtered rings, with supplied iPRA proofs that it is a ring homomorphism,
sends $F^1B$ into $F^1C$, and sends the polynomial system
$P$ to $Q$.  Let $b$ be the root produced in $B$.  Suppose a second valid
lifting instance in $C$ has initial approximation $c$ satisfying
$\rho(a)-c\in F^1C$.  Then $\rho(b)$ equals its produced root.
\end{corollary}

\begin{proof}
Finite polynomial evaluation commutes with the homomorphism, so
$Q(\rho(b))=\rho(P(b))=0$.  Moreover $\rho(b)-c\in F^1C$.
The uniqueness part of Theorem~\ref{thm:filtered-lifting}, applied only in
$C$, identifies this root with the one produced there.  Different lifts
of the initial approximation or of the inverse matrix therefore give the
same answer whenever they satisfy their local hypotheses.
\end{proof}

In a tower of finite quotients, construct the root independently at each
requested precision.  The corollary proves that truncating one answer gives
the answer at lower precision.  This produces a coherent family by one
fixed PR formula.  There is no step of choosing one element from each of
infinitely many nonempty sets.  Inside iPRA, coherence is an arithmetic
statement about any two requested numerical precisions.

\subsection{Ordinary jets and simultaneous implicit functions}

Take
\[
 B_N=R[T]/(T^N),\qquad
 F^kB_N=\{f:[T^j]f=0\text{ for every }j<\min(k,N)\},
                                                               \tag{4}
\]
where $R$ is any fixed PR-presented commutative ring.  An element is a
length-$N$ coefficient list.  Addition is entrywise and multiplication is
the truncated convolution
\[
 [T^n](fg)=\sum_{j=0}^n [T^j]f\,[T^{n-j}]g\quad(n<N).
\]
All filtration laws follow by bounded sums.  For the product law, if
$n<i+j$, every summand has either its first index below $i$ or its second
index below $j$; that factor is zero.  This is a bounded decidable case
split.  Truncation is a ring homomorphism.

\begin{corollary}[Finite simultaneous formal implicit functions]
\label{cor:formal-implicit}
Let $P(T,Y_1,\ldots,Y_m)$ be a polynomial system over $R$.  Supply
$a\in R^m$ and $V\in R^{m\times m}$ such that
\[
 P(0,a)=0,\qquad VJ_Y P(0,a)=J_Y P(0,a)V=I.
\]
For each $N\ge1$, a uniform primitive recursive construction returns the
unique length-$N$ vector of coefficient lists $Y^{[N]}$ with constant
coefficient $a$ and
\[
 P(T,Y^{[N]})=0\pmod {T^N}.
\]
The lists agree under truncation.  The result is valid in arbitrary
characteristic and uses no division by coefficient indices.
\end{corollary}

\begin{proof}
Embed $a,V$ as constant lists and apply Theorem~\ref{thm:filtered-lifting}
in (4).  The hypotheses become its three membership tests at level one.
Apply Corollary~\ref{cor:filtered-coherence} to the truncation maps.
\end{proof}

For a general system $Y=a+T G(T,Y)$, the Jacobian at $T=0$ is the identity.
Thus the theorem handles mutually recursive algebraic coefficient
definitions with no inversion input at all.  The single equation
$Y=1+TY^2$ gives the Catalan recurrence.  The pair
\[
 Y=1+TZ^2,\qquad Z=1+TYZ
\]
illustrates that simultaneous recursion is treated by the same program
and proof; no ordering of the unknown functions has to be chosen.

For inversion, if $f(0)=c$ and $vc=1$, take $P(T,Y)=f(T)Y-1$ and
$a=U=v$.  Put $q=1-vf\in F^1$.  Iteration gives explicitly
\[
 Y^{[N]}=v\sum_{j=0}^{N-1}q^j,
 \qquad fY^{[N]}=(1-q)\sum_{j=0}^{N-1}q^j=1-q^N=1.
\]
This proves the usual inverse-series construction from a finite geometric
identity.  For $P=Y^2-(1+T)$, $a=1$ and $U=1/2$ give the formal square-root
coefficients over a ring where $2$ has a supplied inverse.  Over a ring
where it has none, this particular instance does not satisfy the theorem's
hypotheses; the interface does not hide the obstruction.

These are statements about exact finite jets.  They do not imply an
analytic convergence radius, a real root, or completeness of all singular
branches.  The Puiseux packet's ramification and exact coefficient-field
backend, and its majorant estimates, remain separate mathematics.
Likewise, this theorem does not reprove the Lagrange coefficient formulas;
it supplies a common construction and uniqueness layer beneath them.

\subsection{Modular Hensel lifting, including composite bases}

Let $q\ge2$ and
\[
 B_N=\mathbb Z/q^N\mathbb Z,\qquad
 F^kB_N=\{x:q^{\min(k,N)}\mid x\},            \tag{5}
\]
using the canonical representative $0\le x<q^N$ in the membership test.
Changing a representative by a multiple of $q^N$ preserves this test.
Divisibility, modular operations, powers, and their finite arithmetic
laws are primitive recursive.  The filtration product law follows by
multiplying two divisibility witnesses; if $i+j\ge N$, the product
is zero in $B_N$.

\begin{corollary}[Finite modular lifting]
\label{cor:modular-hensel}
Let $P\in\mathbb Z[Y_1,\ldots,Y_m]^m$ and supply an integer vector $a$
and matrix $U$ satisfying
\[
 P(a)\equiv0\pmod q,\qquad
 J(a)U\equiv UJ(a)\equiv I\pmod q.
\]
For each $N\ge1$, the recurrence
\[
 b_0=a\bmod q^N,\qquad
 b_{r+1}=b_r-UP(b_r)\bmod q^N
\]
returns after $N-1$ updates the unique root modulo $q^N$ in the residue
class of $a$ modulo $q$.  These outputs are compatible under reduction.
No primality assumption on $q$ is required.
\end{corollary}

\begin{proof}
Apply Theorem~\ref{thm:filtered-lifting} in (5), and use the natural
reduction homomorphism for compatibility.  All inverse information is
contained in the displayed finite congruence, which can be checked before
the iteration.  No claim is made that every root modulo $q$ has such a
Jacobian inverse.
\end{proof}

For example, take $q=4$, $P(X)=X^3-X-2$, $a=2$, $U=3$, and $N=3$.
Here $P(2)=4$ and $3P'(2)=33\equiv1\pmod4$.  Modulo $64$ the two
updates are $2\mapsto54\mapsto22$.  Directly,
$22^3-22-2=10624=166\cdot64$.  The theorem proves uniqueness among the
residues congruent to $2$ modulo $4$, not uniqueness among all residues.

\subsection{Hurwitz coefficients and differential equations without division}

There is a second coefficient algebra particularly suitable for integral
combinatorics.  A length-$N$ Hurwitz list has multiplication
\[
 (a\star b)_n=\sum_{j=0}^n\binom nj a_jb_{n-j}\qquad(n<N),       \tag{6}
\]
with identity $(1,0,\ldots,0)$.  Binomial coefficients are defined by
Pascal recursion.  This algebra is valid over any commutative ring and
does not require factorial denominators.  Commutativity follows from
$\binom nj=\binom n{n-j}$.  Associativity follows by collecting the terms
$a_i b_j c_k$ with $i+j+k=n$ and using
\[
 \binom n{i+j}\binom{i+j}i
   =\binom ni\binom{n-i}j.
\]
This identity counts the same ordered division of $n$ positions into
three labeled parts; alternatively it follows by Pascal induction.  Both
the counting sets and the induction invariants are finite.  The filtration
by vanishing initial coefficients satisfies exactly the laws used above,
so Theorem~\ref{thm:filtered-lifting} applies to this algebra too.

Differentiation is the shift $(Da)_n=a_{n+1}$.  Its finite-precision type
must be respected:
\[
 D:\operatorname{Hurwitz}_{N+1}(R)\longrightarrow
       \operatorname{Hurwitz}_{N}(R).                         \tag{7}
\]
It is not an endomorphism of a fixed truncation. Explicitly, with
$a,b$ of length $N+1$, the typed identity is
\[
 D(a\star_{N+1}b)=Da\star_N\pi_N(b)+\pi_N(a)\star_N Db,
\]
where $\pi_N$ truncates to length $N$. Pascal's identity gives
$D(a\star b)=Da\star b+a\star Db$ at every coefficient for which the
needed input coefficients have been supplied.  Indeed, split
$\binom{n+1}j$ into $\binom nj+\binom n{j-1}$ and reindex the two finite
sums.  The demand for one extra coefficient is part of the contract.

\begin{theorem}[Triangular Hurwitz differential systems]
\label{thm:hurwitz-ode}
Fix a PR-presented commutative ring $R$, an initial vector $c\in R^m$,
and a polynomial system $G(T,Y)$ whose coefficient lists are supplied at
the requested finite precision.  There is a unique coefficient prefix
through degree $N$ satisfying
\[
 Y_0=c,\qquad DY=G(T,Y)\quad\text{through degree }N-1.          \tag{8}
\]
The prefix is produced by a PR recurrence, uniformly in $N$, the finite
polynomial code, and the finite coefficient input.  Compatible input
prefixes give compatible output prefixes.  No integer division is used.
\end{theorem}

Here the time variable $T$ is the Hurwitz list $(0,1,0,\ldots)$.
If the coefficients of $G$ are themselves time-dependent lists, they
are additional finite inputs, with their prefixes supplied through
degree $N-1$. Evaluation uses the Hurwitz product throughout. In
particular $T^{\star r}$ has coefficient $r!$ in position $r$;
one does not identify it with the ordinary monomial basis in positive
characteristic.

\begin{proof}
Assume coefficients $Y_0,\ldots,Y_n$ are already stored.  Form each
coefficient of the finite expression $G(T,Y)$ through degree $n$, using
(6) and the stored lists, and set
\[
 Y_{n+1}=[G(T,Y)]_n.                                           \tag{9}
\]
This expression only reads coefficients through degree $n$: at a product,
every convolution index is at most $n$, and addition and constants have
the same property.  Thus (9) is a triangular recurrence, not an implicit
unbounded search.  A single recursion accumulates the finite vector of
prefixes.  Each step makes one new instance of (8) true and preserves all
earlier instances because the earlier coefficients are unchanged.

For uniqueness, suppose another finite prefix satisfies (8).  Its zeroth
coefficient is $c$.  If its coefficients agree through $n$, finite
evaluation of $G$ agrees at coefficient $n$, so (9) forces agreement at
$n+1$.  Bounded-prefix induction proves equality of the whole requested
prefix.  The same argument proves compatibility upon shortening the
coefficient input.  The invariants are conjunctions over $j\le n$ and
$i<m$ of coded ring equalities, so they have quantifier-free Boolean forms.
\end{proof}

For $Y'=A\star Y$, $Y_0=1$, formula (9) becomes
\[
 Y_{n+1}=\sum_{j=0}^n\binom nj A_jY_{n-j}.                      \tag{10}
\]
Taking $A_j=1$ gives the Bell recurrence, while $Y'=Y\star Y$, $Y_0=1$,
gives $Y_n=n!$: every one of the $n+1$ summands
$\binom nj j!(n-j)!$ equals $n!$.
The simultaneous system $Y'=Z$, $Z'=-Y$ gives the usual sine and cosine
coefficient patterns from their initial values without dividing by any
factorial.

Over a ring containing the rationals, the map
$a\mapsto\sum a_n T^n/n!$ identifies (6) with ordinary multiplication
of exponential generating functions.  Over $\mathbb Z$ or a ring of
positive characteristic, that notation is only a motivation; the
division-free algebra (6) is the definition.  Analytic convergence and
tail estimates are additional contracts.  This recovers the finite
recurrence layer of the Bell packet and generalizes it to arbitrary
polynomial differential systems, without claiming its analytic asymptotic
results anew.

\subsection{Lean definitions and completion obligations}

The finite data layer should define the following objects before invoking
any host-language theorem about rings:
\begin{enumerate}
\item \texttt{RingCode}: numerical validity and equality tests; codes for
zero, one, addition, negation and multiplication; explicit PR syntax for
each of these fixed operations.
\item \texttt{FilteredCode}: a precision parameter, the Boolean predicate
\texttt{member k x}, and the supplied object-iPRA derivations of the
filtration laws above.  \texttt{VecMember} and \texttt{MatMember} are
bounded Boolean conjunctions.
\item \texttt{PolyCode}: a list of coefficient/exponent-vector pairs;
\texttt{eval}, \texttt{partial}, \texttt{jacobian},
\texttt{matVec}, and \texttt{matMul} are explicit PR folds.  A
\texttt{GoodPoly} predicate checks every arity and carrier condition.
\item \texttt{liftState P U a r}: the recursion (L), with a default value
on malformed input; \texttt{LiftHyp} is the bounded conjunction (H).
\item \texttt{JetCode N} and \texttt{ModCode q N}: the concrete
implementations (4) and (5), each with proofs of the abstract interface.
\texttt{HurwitzCode N} replaces ordinary convolution by (6).
\end{enumerate}

The main mathematical targets are the following explicit propositions,
with all variables understood as numerical codes:
\begin{align*}
 &\operatorname{LiftHyp}(N,P,U,a)\ \longrightarrow\
   \operatorname{Eval}(P,\operatorname{liftState}(P,U,a,N-1))=0,\\
 &\operatorname{LiftHyp}(N,P,U,a)\wedge P(x)=0\wedge x-a\in F^1
      \ \longrightarrow\
   x=\operatorname{liftState}(P,U,a,N-1).
\end{align*}
Their deep endpoints are terms of
\texttt{Prf PRAAx [] (allClose LiftExistence)} and
\texttt{Prf PRAAx [] (allClose LiftUniqueness)}, after the supplied
presentation laws have been discharged or retained explicitly as premises.
They are not statements that an arbitrary host ring can be automatically
turned into an arithmetic presentation.

The required intermediate lemmas are, in order: finite sum and matrix
filtration closure; the uniform finite Taylor expansion in
Lemma~\ref{lem:filtered-taylor}; residual precision improvement (1);
the Boolean invariant (3); difference improvement (2); homomorphism
commutation with finite evaluation; and the two concrete implementations.
For Hurwitz equations add Pascal's identity, the triple-binomial identity,
coefficient causality for \texttt{eval}, and the bounded-prefix invariant
for (9).  Each item has its numerical definition and its exact induction
formula specified above.  The general theorem and these implementations
are supplied here as human proofs and a detailed formalization plan;
they are not reported as already compiled deep iPRA derivations.

\section{Weighted polynomial norms: an arithmetic zero-count engine}
\label{sec:weighted-norm}

The previous framework packet proves a sharp zero count for equations
$Y^2=f(X)$ with $\deg f$ odd \cite{repoEmbedding}. The construction below extends that
mechanism to an entire class of plane curves, including equations with
mixed $X,Y$ terms.  It also gives an exact dimension formula for a family
of finite spaces of functions.  The arguments require only polynomial
arithmetic, determinants, exact rational comparisons, and induction on
bounded lists.  They use neither a divisor theorem nor a choice of
algebraic branches.

\subsection{Input data and the exact norm degree}

Fix coprime positive integers $m,d$ with $m\geq2$, and a nonzero rational
$c$.  Give $X$ weight $m$ and $Y$ weight $d$.  Consider
\begin{equation}\label{eq:weighted-curve}
 F(X,Y)=Y^m-cX^d+
       \sum_{mi+dj<md} f_{ij}X^iY^j,
 \qquad f_{ij}\in\mathbb Q.
\end{equation}
The sum is a finite coefficient table; missing coefficients mean zero.
All its $Y$ exponents are less than $m$.  Consequently $F$ is monic of
$Y$ degree $m$ and the quotient
\[
 \mathcal A=\mathbb Q[X,Y]/(F)
\]
is a free $\mathbb Q[X]$ module with the explicit basis
$1,Y,\ldots,Y^{m-1}$.  Here freeness is proved directly by division by a
monic polynomial: subtract the leading multiple of $F$ until the $Y$
degree is below $m$; a nonzero multiple of the monic $F$ cannot itself
have $Y$ degree below $m$.

A reduced polynomial is an expression
\[
 G(X,Y)=\sum_{j<m} A_j(X)Y^j.
\]
Suppose $G\ne0$ and let
\begin{equation}\label{eq:weighted-size}
 W(G)=\max\{mi+dj:[X^iY^j]G\ne0\}.
\end{equation}
Write $M_G(X)$ for the $m$ by $m$ matrix of multiplication by $G$ on
$\mathcal A$ in the displayed basis, and put
\[
 N_G(X)=\det M_G(X).
\]
These are explicit polynomials.  To compute column $j$, reduce $GY^j$
by the equation $Y^m=cX^d-\sum f_{ij}X^iY^j$.  Its starting $Y$ degree
is at most $2m-2$; each reduction lowers the largest remaining $Y$
exponent.  The determinant can be defined by its finite permutation sum.
This definition is useful even when a faster elimination algorithm is
used by an implementation.

\begin{theorem}[Exact weighted norm degree]\label{thm:weighted-norm-degree}
There is a unique monomial $aX^kY^j$ of $G$ of maximum weight, with
$a\ne0$ and $0\le j<m$.  Its norm satisfies
\[
 \deg N_G=W(G)=mk+dj,
 \qquad
 \operatorname{lc}(N_G)
   =a^m\bigl((-1)^{m-1}c\bigr)^j\ne0.
\]
In particular the matrix determinant is nonzero, including when the curve
has singular affine points.
\end{theorem}

\begin{proof}
If $mi+dj=mi'+dj'$ with $0\le j,j'<m$, then $m$ divides
$d(j-j')$.  Coprimality implies that $m$ divides $j-j'$, and its absolute
value is less than $m$.  Thus $j=j'$ and $i=i'$.  Distinct reduced
monomials therefore have distinct weights.  This proves the uniqueness
of the maximum.

The leading determinant can be calculated without guessing which of its
summands cancel.  Introduce an indeterminate $t$ and define
\begin{align*}
 F_t(Z)&=t^{md}F(t^{-m},t^{-d}Z),\\
 G_t(Z)&=t^{W(G)}G(t^{-m},t^{-d}Z).
\end{align*}
The weight inequalities show that both expressions lie in
$\mathbb Q[t,Z]$: a monomial of weight $v$ is multiplied by
$t^{md-v}$ or by $t^{W(G)-v}$.  Moreover $F_t$ is monic of degree $m$
in $Z$, and specialization at $t=0$ gives
\[
 F_0(Z)=Z^m-c,\qquad G_0(Z)=aZ^j.
\]
Let $D(t)$ be the determinant of multiplication by $G_t$ on
$\mathbb Q[t,Z]/(F_t)$ in its monic-remainder basis.  This is a polynomial
in $t$.  Over the Laurent polynomial ring, the change of basis
$Y=t^{-d}Z$ is invertible, so conjugating the multiplication matrices and
factoring the scalar $t^{-W(G)}$ out of their $m$ columns gives
\begin{equation}\label{eq:weighted-rescaling}
 D(t)=t^{mW(G)}N_G(t^{-m}).
\end{equation}
This identity can alternatively be verified directly on the $m$ basis
vectors; no semantic assertion about a field of algebraic functions is
needed.

At $t=0$, multiplication by $Z$ in
$\mathbb Q[Z]/(Z^m-c)$ has the cyclic companion matrix whose determinant
is $(-1)^{m-1}c$.  Multiplication by $aZ^j$ is $a$ times the $j$th power
of that matrix.  Determinant multiplicativity therefore yields
\[
 D(0)=a^m\bigl((-1)^{m-1}c\bigr)^j\ne0.
\]
Write $N_G(X)=\sum_{v=0}^{q}n_vX^v$ with $n_q\ne0$, if the
polynomial is nonzero.  Equation~\eqref{eq:weighted-rescaling} first
shows that it cannot be zero.  Its right side has terms
$n_vt^{m(W(G)-v)}$.  Distinct $v$ have distinct exponents.  Since the
left side has no negative powers, $q\le W(G)$.  Since its constant
coefficient is nonzero, $q=W(G)$, and that coefficient is the displayed
leading coefficient of $N_G$.  This proves all assertions.
\end{proof}

\paragraph{An immediate algebraic benefit.}
The quotient $\mathcal A$ has no zero divisors.  Indeed, for a nonzero
reduced $G$, the adjugate identity says
$\operatorname{adj}(M_G)M_G=N_G I$.  If $GH=0$, it follows that every
coefficient of the reduced $H$ is annihilated by the nonzero polynomial
$N_G$.  Rational polynomial arithmetic has no zero divisors, so $H=0$.
Thus an exact function-field interface can be built from pairs of
reduced polynomials with a nonzero denominator.  Equality is tested by
reducing cross products, and operations are finite polynomial
computations.  This observation supplies an algebraic interface; it does
not claim a general implementation of divisors or Riemann--Roch spaces.

\subsection{Counting points with repeated first coordinates}

A determinant root bound by itself counts distinct $X$ coordinates, not
points.  The following argument records the necessary multiplicities.

\begin{lemma}[A fiber contributes its full cardinality]\label{lem:weighted-fiber}
Fix $u\in\mathbb Q$.  If the distinct rationals
$v_1,\ldots,v_r$ satisfy
\[
 F(u,v_i)=G(u,v_i)=0\qquad(1\le i\le r),
\]
then $(X-u)^r$ divides $N_G(X)$.
\end{lemma}

\begin{proof}
The polynomial $F(u,Y)$ is monic of degree $m$, so the elementary
polynomial root bound gives $r\le m$.  Define row vectors
\[
 e(v_i)=(1,v_i,\ldots,v_i^{m-1}).
\]
Evaluation is compatible with reduction modulo $F(u,Y)$, because
$F(u,v_i)=0$.  Applying the evaluation row to each column of the
multiplication matrix therefore gives
\[
 e(v_i)M_G(u)=G(u,v_i)e(v_i)=0.
\]
Form an $m$ by $m$ rational matrix $U$ whose first $r$ rows are these
vectors and whose remaining rows are the standard basis rows with
indices $r,r+1,\ldots,m-1$.  Its determinant is the determinant of the
first $r$ columns of its first $r$ rows, namely
\[
 \det U=\prod_{1\le i<j\le r}(v_j-v_i)\ne0.
\]
The empty product is one.  Thus no unspecified basis extension is used.

The first $r$ rows of $UM_G(X)$ vanish at $X=u$.  The factor theorem
makes every entry of those rows divisible by $X-u$.  Factoring this
factor from each of the $r$ rows shows that $(X-u)^r$ divides
$\det(UM_G)=(\det U)N_G$.  Since $\det U$ is a nonzero rational,
it is invertible, and the same divisibility holds for $N_G$.
\end{proof}

\begin{theorem}[Sharp arithmetic zero count]\label{thm:weighted-zero-count}
Every finite list of pairwise distinct rational points satisfying
$F(x,y)=G(x,y)=0$ has length at most $W(G)$.
\end{theorem}

\begin{proof}
Sort the given finite list by its rational first coordinate, and let
$u_1,\ldots,u_s$ be the distinct first coordinates, with fiber sizes
$r_1,\ldots,r_s$.  Inside each fiber the second coordinates are distinct,
because the points are distinct.  Lemma~\ref{lem:weighted-fiber} gives
$(X-u_i)^{r_i}\mid N_G$ for every $i$.

The powers belonging to different $u_i$ are coprime.  For completeness,
this can be checked without factoring a general polynomial.  If
$P=(X-u)^r Q$ and $(X-v)^s\mid P$ with $u\ne v$, write the identity
in powers of $T=X-v$.  The first factor has the nonzero constant
coefficient $(v-u)^r$.  If the first $s$ coefficients of $P(v+T)$
vanish, the constant coefficient equation first forces $Q(v)=0$;
the coefficient equation at $T^k$ then forces the $k$th coefficient
of $Q(v+T)$ to vanish once the earlier ones vanish.  Induction on
$k<s$ proves $(X-v)^s\mid Q$.  Repeating this removal for the finite
list of fibers proves
\[
 \prod_{i=1}^s(X-u_i)^{r_i}\mid N_G.
\]
The product has degree $\sum_i r_i$, the length of the input list.
Theorem~\ref{thm:weighted-norm-degree} says that the nonzero polynomial
$N_G$ has degree $W(G)$, so $\sum_i r_i\le W(G)$.
\end{proof}

The proof handles ramification and singular fibers automatically.  It does
not assume that $F(u,Y)$ has $m$ distinct roots, that it splits over the
rationals, or that the roots can be selected continuously.

\subsection{Finite spaces and an elementary dimension formula}

For a numerical budget $L\ge0$ define the explicit monomial list
\[
 \mathcal B_L=\{X^iY^j:0\le j<m,\ i\ge0,\ mi+dj\le L\}.
\]
Every entry has $i\le L$ and $j<m$, so a rectangular bounded scan
constructs it.  Its elements are linearly independent in $\mathcal A$
by the monic-remainder uniqueness proved above.  Hence the dimension of
their span equals the directly computable integer
\begin{equation}\label{eq:weighted-dimension}
 n(L)=\sum_{j=0}^{m-1}
       \max\left(0,1+\left\lfloor\frac{L-dj}{m}\right\rfloor\right).
\end{equation}
In this displayed formula the floor is the signed integer floor.  A
natural-number program uses the equivalent branch: contribute zero if
$L<dj$, and otherwise contribute $1+\lfloor(L-dj)/m\rfloor$.

\begin{proposition}[Exact eventual dimension]\label{prop:weighted-dimension}
Put $g_0=(m-1)(d-1)/2$.  This is an integer.  For every
$L\ge\max(0,md-m-d)$,
\[
 n(L)=L+1-g_0.
\]
\end{proposition}

\begin{proof}
Since $m,d$ are coprime, they cannot both be even, so
$(m-1)(d-1)$ is even.  Under the stated lower bound on $L$,
$L-dj\ge -m$ for all $j<m$.  Therefore each summand
$1+\lfloor(L-dj)/m\rfloor$ is nonnegative and the maximum with zero
can be removed.

The residues $r_j$ of $L-dj$ modulo $m$ are a permutation of
$0,1,\ldots,m-1$, by coprimality.  Write
$L-dj=mq_j+r_j$, $0\le r_j<m$.  Summing these $m$ equations gives
\[
 m\sum_{j<m}q_j
   =mL-\frac{dm(m-1)}2-\frac{m(m-1)}2.
\]
Adding $m$ to $\sum_jq_j$ proves the claimed formula.
Every calculation is an identity of finite sums; no geometric dimension
or genus theorem is invoked.
\end{proof}

The number $g_0$ agrees with familiar geometric formulas in suitable
smooth cases, but no identification with the genus is needed here.  The
proposition proves exactly the numerical dimension assertion stated.

For a rational point $P=(x,y)$ of the curve, let
$\Phi_L(P)$ be the vector of the values of the monomials in
$\mathcal B_L$, in the scan order.  Because $1\in\mathcal B_L$,
this vector is nonzero.

\begin{corollary}[Uniform bound over a proper subspace]\label{cor:weighted-subspace}
Let $V$ be a proper rational linear subspace of
$\mathbb Q^{n(L)}$, specified by a finite rational basis or matrix.
Any finite list of distinct curve points whose evaluation vectors lie
in $V$ has length at most $L$.
\end{corollary}

\begin{proof}
Exact Gaussian elimination produces a nonzero rational row vector
$\lambda$ annihilating $V$: reduce the finite system of linear
equations for the annihilator, choose a free variable, set it to one,
and solve backwards with witnessed nonzero pivots.  Properness means
that a free variable exists.  Define
\[
 G_\lambda(X,Y)=\sum_{b\in\mathcal B_L}\lambda_b b(X,Y).
\]
It is a nonzero reduced polynomial, because the monomials are distinct
and $\lambda\ne0$, and it has weight at most $L$.  All the given points
are zeros of $G_\lambda$.  Apply
Theorem~\ref{thm:weighted-zero-count}.
\end{proof}

\begin{corollary}[The remaining approximation contract is isolated]\label{cor:weighted-cover}
Let $S$ be a decidable class of points on the curve.  Suppose every
finite list of distinct points in $S$ has its $\Phi_L$ vectors covered
by at most $M$ supplied proper rational subspaces.  Then every such list
has length at most $ML$.
\end{corollary}

\begin{proof}
For each point scan the finite list of subspaces and assign it to the
first one containing its vector.  Subspace membership is decidable by
finite exact elimination, and the assumed cover proves that the scan
succeeds.  The points assigned to one subspace form a distinct sublist,
so the preceding corollary bounds that sublist by $L$.  The original
length is the sum of at most $M$ such lengths, hence at most $ML$.
\end{proof}

This resolves the zero-count side of a finite-cover argument for every
curve of the displayed weighted form.  Producing or proving the stated
approximation cover is a separate dependency.  Neither the norm theorem
nor the dimension formula produces that cover.  The conclusion is a
cardinality bound on arbitrary supplied lists, not a height bound or a
complete enumeration of all points.

\subsection{Examples and sharpness}

Take $m=3$, $d=4$ and
\[
 F(X,Y)=Y^3-X(X-1)(X-2)(X-3).
\]
All terms other than $Y^3-X^4$ have weight below $12$.  For $G=Y$,
$W(G)=4$ and
\[
 N_G(X)=X(X-1)(X-2)(X-3).
\]
The four points $(0,0),(1,0),(2,0),(3,0)$ are distinct common zeros.
Thus the zero-count bound is attained.  Here
$\mathcal B_4=(1,X,Y)$ in a suitable order, and the proper subspace
with last coordinate zero has exactly four points in its inverse image.

The same method includes mixed equations such as
\[
 Y^3=X^4+XY+X+1.
\]
The mixed monomial $XY$ has weight $7<12$.  Every nonzero linear form
$\alpha+\beta X+\gamma Y$ has at most four rational zeros on this
curve, and every nonzero combination of
$1,X,Y,X^2,XY,Y^2$ has at most eight.  These assertions follow from the
weight bounds, with no test of smoothness and no root isolation.
For $m=3,d=4$ the dimension formula gives $n(L)=L-2$ for $L\ge5$;
for example $\mathcal B_8$ has six elements.

\subsection{A broader finite certificate when coprimality fails}

The rescaling argument has a useful form that does not assume
$\gcd(m,d)=1$.  Allow any monic $F$ of $Y$ degree $m$ whose monomials
have weight at most $md$, and any reduced nonzero $G$ of maximum
weight $W$.  Form $F_t,G_t$ by the same formulas, and compute
\[
 C(F,G)=\det\left(\text{multiplication by }G_0
                 \text{ on }\mathbb Q[Z]/(F_0)\right).
\]
Here $F_0$ is still monic of degree $m$.  If the exact rational test
$C(F,G)\ne0$ succeeds, the identical proof gives
\[
 \deg N_G=W,\qquad
 \#\{\text{distinct supplied common zeros}\}\le W.
\]
Thus coprimality is a simple sufficient condition making this test
succeed for every nonzero reduced $G$ in our principal family.  For
more general weighted curves the finite initial determinant is a usable
certificate on individual linear forms.  Failure of the test means only
that this leading-degree certificate did not apply; it does not decide
the number of zeros or the irreducibility of the original curve.

\subsection{The iPRA translation and a Lean development path}

The theorem is uniform in finite rational coefficient tables and in the
natural numbers $m,d,L$.  It must not be represented by a hypothetical
universal evaluator for arbitrary primitive recursive programs.  Instead
fix the finitely many algorithms below and represent each by its own
primitive recursive definition.

\begin{enumerate}
\item Encode signed integers, normalized fractions, coefficient lists,
matrices and point lists.  Implement rational comparison and equality by
cross multiplication with positive denominators.  Prove the totalized
operations correct on the decidable valid-code predicate.
\item Define rational polynomial addition, multiplication, coefficient
lookup, degree, evaluation, monic division in $Y$, and the finite
permutation determinant.  Prove their coefficient identities with
bounded folds.  A determinant specialization lemma and determinant
multiplicativity are explicit finite identities, not imported reflection
principles.
\item Construct $M_G$ column by column.  Prove the basis evaluation
identity $e(v)M_G(u)=G(u,v)e(v)$ from the monic-remainder identity,
under the two numerical polynomial-evaluation hypotheses.
\item Formalize the weight uniqueness lemma, the nonnegative exponent
checks for $F_t,G_t$, and the diagonal change-of-basis identity.  Negative
powers can be encoded as a finite Laurent coefficient table with a signed
starting exponent.  Alternatively multiply the identity by a computed
power of $t$ and work exclusively with ordinary polynomial coefficients.
\item Prove the explicit cyclic-matrix determinant, obtaining the nonzero
constant term.  Compare coefficients in the rescaling identity to derive
the norm degree and leading coefficient.
\item Prove the Vandermonde determinant and the fiber-divisibility lemma
using the explicit matrix $U$.  Prove the successive factor-removal
lemma by the triangular coefficient argument.  Apply it to the sorted
input point list to obtain the zero-count inequality.
\item Prove the finite count of monomials and the residue-permutation
identity.  Define a matrix presentation of a proper subspace, implement
rational row elimination with at most the matrix dimensions many pivot
steps, and build the annihilating row.  This gives the subspace and
finite-cover corollaries.
\end{enumerate}

For the object-language proof, every loop invariant can be expressed as
a Boolean equation about a bounded prefix: coefficient identities,
nonzero pivots, the residual factorization, and the current degree bound.
Parameters such as a point, a coefficient table or an input list remain
free during the induction.  Universal quantification over those
parameters is introduced afterwards.  No induction over an unbounded
statement about all points or all polynomials is required.

The accompanying \path{WeightedNormTargets.lean} defines exact rational
tables and states the propositions \path{NormDegreeGoal},
\path{FiberPowerGoal}, \path{CommonZerosGoal},
\path{WeightedBasisCountGoal}, and \path{EventualDimensionGoal}.
The proper-subspace corollary additionally uses the finite row-elimination
construction specified above. An ambient Lean proof of these
statements verifies their mathematical algebra.  To establish their
asserted iPRA derivability, the coefficient and loop proofs still have
to be translated into the repository's embedded proof judgment, and
then transported through its checked theory translation.  This section
specifies that route and gives the human proof; it does not assert that
those embedded derivation objects have already been constructed.

\section{Lean formalization: definitions, verified core, and completion path}
\label{sec:formalization}

\subsection{The endpoint is an embedded proof}

The included standalone kernel reproduces three current repository files:
\texttt{Syntax.lean}, \texttt{Deduction.lean}, and \texttt{PR.lean}.
The source bodies were checked against the pinned commit. Original imports
are replaced by \texttt{import Std} and a natural-number notation; the
declarations and proofs are preserved. The archive retains the original
source files, their provenance and the repository license.

The key interfaces, in namespace \texttt{LeanProofs.ConciseIPRA}, are:
\begin{lstlisting}
PR k                 -- a fixed primitive recursive function symbol
PTm n                -- a term with n free variables
Fm PTm n             -- a first-order arithmetic formula
Prf PRAAx Gamma phi  -- an object-iPRA derivation
PRA.qfInd            -- quantifier-free object induction
\end{lstlisting}
Here \texttt{PR} is an inductive syntax with constructors for zero,
successor, projection, composition and primitive recursion. A constructor
of type \texttt{PR k} supplies a fixed program symbol, not an interpreter
for arbitrary programs. \texttt{Prf} is the intuitionistic natural
deduction relation with equality and the stated axiom scheme.

The current repository proves the exact transport statement
\begin{lstlisting}
pra_iff_native :
  (Gamma |- [PRAAx] phi) <->
  NPrf (Gamma.map NFm.ofPRA) (NFm.ofPRA phi)
\end{lstlisting}
where the displayed ASCII turnstile is only explanatory notation for
the repository's scoped judgment \cite{repoBridge}. The forward direction
is applied to an already constructed arithmetic derivation. It does not
accept a host theorem about naturals as a replacement. The standalone
build does not rebuild the full native equivalence dependency closure.

\subsection{What is checked in the supplied Lean core}

\begin{longtable}{P{.27\linewidth}P{.43\linewidth}P{.21\linewidth}}
\toprule
Module & Contents & Logical status\\
\midrule\endhead
{\footnotesize\path{ObjectKernel.lean}} & Pinned scoped syntax, substitution,
intuitionistic deduction and ordinary iPRA axioms. & Reused foundations.\\[3pt]
{\footnotesize\path{RankedDescent.lean}} & Run validity, rank budget, termination in the
initial rank, invariant transport and a contraction-to-descent lemma.
& Proved host theorems.\\[3pt]
{\footnotesize\path{CoherentTower.lean}} & Iteration of specified local lifts,
admissibility, adjacent coherence, base projection and uniqueness over
the seed. & Proved host theorems.\\[3pt]
{\footnotesize\path{ObjectComposition.lean}} & Witness-contract composition,
specialization, existential introduction, a QF-induction interface and
an instantiated identity-iteration proof. & Actual object-proof
constructions.\\[3pt]
{\footnotesize\path{DescentTargets.lean}} & Explicit PR codes for arithmetic and bounded
iteration; exact object formulas and the remaining descent compiler
goal. & Definitions; the compiler goal is unproved.\\[3pt]
{\footnotesize\path{ApplicationTargets.lean}} & Polynomial systems, Jacobians, ring and
filtration laws, frozen iteration, jets, residues and projections.
& Compiled definitions; general lifting and presentation goals unproved.\\[3pt]
{\footnotesize\path{WeightedNormTargets.lean}} & Rational polynomials, monic reduction,
multiplication matrices, determinants, weights and monomial counts.
& Compiled definitions; general norm and counting goals unproved.\\[3pt]
{\footnotesize\path{HurwitzTargets.lean}} & Binomial products, typed derivative demand,
and a scalar rational polynomial-ODE prefix producer.
& Compiled definitions; general algebra and ODE goals unproved.\\[3pt]
{\footnotesize\path{ApplicationChecks.lean}} & A modular root and a finite basis count;
exact rational jet and norm computations. & Two finite kernel theorems;
rational diagnostics are executable checks.\\[3pt]
{\footnotesize\path{AxiomAudit.lean}} & Printed dependencies of the core declarations.
& Reproducible audit.\\
\bottomrule
\end{longtable}

The core build uses Lean 4.32.0 and \texttt{Std}. Its coherence and
uniqueness declarations use no axioms. The other audited core results
use at most Lean's metatheoretic \texttt{propext} and \texttt{Quot.sound};
none uses \texttt{Classical.choice} or \texttt{sorryAx}. Those dependencies
are properties of the host proofs, not extra axioms inserted into
\texttt{PRAAx}. The build and axiom logs list the exact declarations.

All ten modules compile. The application propositions are ordinary
\texttt{Prop} definitions, not axioms or admitted proofs. Their compilation
checks that the definitions and targets are well typed; it does not prove
those targets. The modular example proves by kernel reduction that the
producer for \(X^3-X-2\), with \(q=4\), \(N=3\), seed \(2\) and inverse
\(3\), returns \(22\). The rational diagnostics check the Catalan prefix
\((1,1,2,5,14)\), a mixed-curve norm \(1+X+X^4\), and the Hurwitz
solution prefix \((1,1,2,6,24,120,720,5040)\) for \(Y'=Y^2\).
The concrete ODE target uses a scalar polynomial with constant rational
coefficients. The broader time-dependent, multivariate human theorem
extends the same finite coefficient recurrence; that extension is part
of the documented completion route.

\subsection{A genuine composition theorem}

For formulas \(P(x)\) and \(Q(y,x)\), the actual structure
\texttt{TermContract P Q} contains a fixed term \(f(x)\) and an object
derivation of \(Q(f(x),x)\) from \(P(x)\). A second contract contains
\(g(y,x)\) and a derivation of \(R(g(y,x),y,x)\) from \(Q(y,x)\).
The checked composition operation substitutes \(f(x)\) for \(y\), then
uses cut to produce
\[
        [P(x)]\ \vdash_{\ipra}\ R(g(f(x),x),f(x),x).
\]
Capture avoidance and context substitution are handled by the repository's
scoped syntax. The record contains proof obligations; there is no axiom
asserting that arbitrary preconditions and postconditions can be filled.

The file also defines the explicit PR symbol
\begin{lstlisting}
identityIteration : PR 2 :=
  .prec (.proj (0 : Fin 1)) (.proj (1 : Fin 3))
\end{lstlisting}
and proves in the embedded theory that iterating the identity on a seed
returns the seed. Its base uses \texttt{PRA.precZero} and
\texttt{PRA.proj}; its step uses \texttt{PRA.precSucc}, the accumulator
projection and the induction hypothesis. The outer induction is the
actual \texttt{PRA.qfInd}. This small example checks the full syntactic
route without misrepresenting the scale of the remaining algebra library.

\subsection{Numerical definitions for the application layers}

The following definitions are sufficient; no abstract analytic or
algebraic-geometric object is required as an object-theory primitive.

\begin{description}
\item[Ring and filtration.] A fixed finite tuple of PR symbols provides
\(\Valid(x)\), \(\operatorname{eq}(x,y)\), \(0,1,+,-,\cdot\) and
\(\operatorname{member}(N,k,x)\). The ring and filtration laws are
object derivations or explicit hypotheses. Their parameters include the
chosen precision and, for residues, the base \(q\).
\item[Polynomial system.] A polynomial is a list of pairs
\((c,e)\), where \(e\) is a vector of nonnegative exponents. A bounded
test verifies the vector length and coefficient validity. Evaluation
multiplies powers along each vector and then adds terms; formal partial
derivatives replace \(cX^e\) by
\(ce_jX^{e-\mathbf e_j}\) when \(e_j>0\). The Jacobian is the finite
matrix of evaluations of these partial derivatives.
\item[Jets.] A precision-\(N\) element is a coefficient vector of length
\(N\), with zero-default lookup. Equality compares the entries. Products
are truncated convolution; membership in level \(k\) checks that the
first \(\min(k,N)\) entries are zero. Hurwitz jets insert the natural
factor \(\binom nj\) in each product summand.
\item[Residues.] A precision-\(N\) element is the canonical natural
number below \(q^N\). Operations reduce signed integer expressions
modulo that modulus. The membership test is divisibility by
\(q^{\min(k,N)}\). Matrices and vectors are finite lists of these codes.
\item[Plane polynomials.] The supplied target module uses a finite list
of rational \(X\)-coefficient lists indexed by \(Y\)-degree, with
zero-default lookup. Equivalently, an arithmetic encoding can use sparse
triples \((i,j,a_{ij})\), normalized by adding duplicates and dropping
zeros. The monic remainder is computed by eliminating the largest
\(Y\) exponent. Conversion between the two finite representations is PR.
\item[Multiplication norm.] Construct column \(j<m\) by reducing
\(GY^j\). The matrix is an \(m\)-square array of rational
\(X\)-polynomials. Its determinant is either the signed permutation sum
or a recursively expanded Laplace determinant. Both are finite PR
algorithms; their equivalence can be proved coefficientwise.
\item[Point lists and subspaces.] A point is a pair of rational codes.
The common-zero predicate is evaluation equality. Pairwise distinctness
is the bounded conjunction over pairs of list indices. A subspace is a
finite rational spanning matrix; membership is a rank comparison, and
properness is rank below the ambient dimension. The weight-bounded
monomial list is enumerated by \(j<m\), \(i\le L\) and \(mi+dj\le L\).
\end{description}

Malformed arrays are explicitly excluded by bounded well-formedness tests.
The totalized operations may return defaults on malformed inputs. The
algebraic laws and result statements are always guarded by those tests.

\subsection{From the human lifting proof to a complete Lean proof}

Let \(H(N,P,U,a)\) be the bounded input test consisting of validity,
the residual membership and the two inverse-Jacobian congruences in
Theorem~\ref{thm:filtered-lifting}. Let \(B(N,P,U,a)\) be the state
after \(N-1\) updates. Its correctness target is
\[
 H(N,P,U,a)\to
 \bigl(P(B)=0\ \land\ B-a\in F^1\bigr).
\]
The uniqueness target additionally takes a \emph{fixed numerical candidate}
\(z\) and concludes
\[
 H(N,P,U,a)\land P(z)=0\land z-a\in F^1\quad\to\quad z=B.
\]
Every part is a bounded formula once arrays are encoded. Universal closure
over the input codes can be introduced after its proof.

The finite dependency order is as follows.
\begin{enumerate}
\item Prove the integer, rational and array operations correct on valid
codes. Give their \texttt{PR} syntax and their defining equations.
\item Prove addition and matrix multiplication preserve the indicated
filtration levels. The invariant is a conjunction over the processed
entries and the accumulated sum.
\item Prove polynomial evaluation, partial differentiation and the finite
Taylor identities. Induct on monomial exponents and on the length of
the processed term list; the invariant is a finite coefficient identity
and membership test, not a universally quantified ring identity.
\item Derive the single-step residual improvement equation (1) of the
lifting proof. Freeze \(N\), Booleanize invariant (3), and apply
\texttt{PRA.qfInd} to the explicit PR state recursion.
\item Freeze the two candidate roots and \(N\). Begin the uniqueness
induction at filtration zero, use the seed congruence for the step to
level one, and equation (2) for all later steps. At level \(N\), use
the proved nilpotence test.
\item Prove finite evaluation commutes with restriction. Compose this
with uniqueness to obtain pairwise coherence. This last step needs no
infinite family object.
\item Instantiate the generic contracts for ordinary jets, Hurwitz jets
and residues. For the Hurwitz ODE use a separate prefix recurrence;
prove coefficient causality by induction on the finite polynomial
evaluation, then induct on the stored solution length.
\end{enumerate}

At the end, replace each bounded formula by its fixed Boolean compiler
output and build terms of
\texttt{Prf PRAAx [] (allClose liftingFormula)}. The actual compiler
and the large algebraic derivation are not completed in this archive.
The definitions, local identities and induction formulas above make the
remaining work explicit rather than introducing an uninterpreted
``lifting soundness'' axiom.

\subsection{From the norm proof to a complete Lean proof}

Define \(\operatorname{GoodCurve}(m,d,c,F)\) to check
\(m\ge2\), \(d\ge1\), \(\gcd(m,d)=1\), \(c\ne0\), the two leading
coefficients and every lower-weight restriction. Define
\(\operatorname{ReducedNonzero}(m,G)\) by its finite coefficient tests,
and \(W(m,d,G)\) by bounded maximum over the nonzero terms. The main
object formula is the universal closure of
\[
 \begin{split}
 &\operatorname{GoodCurve}(m,d,c,F)\land
   \operatorname{ReducedNonzero}(m,G)\land
   \operatorname{DistinctCommonZeros}(F,G,L)\\
 &\hspace{35mm}\longrightarrow\quad
   \operatorname{length}(L)\le W(m,d,G).
 \end{split}
\]
Unlike a formula quantifying over a completed function field, this is an
ordinary finite assertion about the supplied polynomial and point codes.

\begin{longtable}{P{.28\linewidth}P{.34\linewidth}P{.29\linewidth}}
\toprule
Lemma to formalize & Finite mechanism & Induction invariant\\
\midrule\endhead
Monic remainder & Eliminate the highest \(Y\)-coefficient using \(F\).
& Original polynomial equals quotient times \(F\) plus residual;
processed high coefficients are zero.\\[3pt]
Norm specialization & Evaluate every entry in the multiplication matrix;
evaluate its finite determinant. & Equality of the accumulated
permutation sums.\\[3pt]
Rescaling identity & Clear a computed common power of \(t\), then compare
finite polynomial coefficients. & The change-of-basis equality for
processed matrix columns.\\[3pt]
Exact norm degree & Compute the cyclic companion determinant and compare
the constant term. & Finite coefficient support and a certified nonzero
top coefficient.\\[3pt]
Fiber multiplicity & Use the explicit Vandermonde rows and factor
\(X-u\) out of their entries. & The row identities and accumulated
factorization.\\[3pt]
Product of fiber factors & Remove distinct-coordinate powers successively.
& \(N_G=C_kQ_k\), with \(C_k\) the processed product, and the finite
Taylor coefficients needed for the next fiber.\\[3pt]
Dimension count & Enumerate the bounded monomials and residues modulo \(m\).
& Prefix counts and the permutation of residue classes.\\[3pt]
Subspace fiber bound & Row-reduce the rational spanning matrix and produce
an annihilator. & Row equivalence, nonzero pivots and the explicit free
variable giving a nonzero kernel vector.\\
\bottomrule
\end{longtable}

For all coefficient identities, the finite support bound is part of the
state. Monic reduction lowers the largest remaining \(Y\)-exponent; in
the principal weighted family it also never increases monomial weight.
The determinant has \(m!\) permutation terms. Rational values may grow
rapidly, but the state is still a PR-coded finite list. There is no need
for a common external bound on every possible rational input value.

An ambient algebra proof can be completed first using ordinary Lean
polynomials and matrices; that would be useful independent validation.
For the intended iPRA theorem, it must then be implemented in the numerical
coefficient library or accompanied by an explicit verified compiler that
constructs the corresponding embedded derivation. The target modules in
this archive make no general reflection assumption.

\subsection{Repository integration and reproducible builds}

The standalone source is intended for inspection and independent builds.
Inside ProveIt, replace the import of \texttt{ObjectKernel} by the current
\texttt{ConciseIPRA.PR} import and build against the pinned toolchain.
Do not import both versions of the same namespace in one build. Once a
new object theorem is completed, apply \texttt{pra\_iff\_native.mp} to
obtain its native form. The provided scripts compile the standalone
sources, run the axiom audit and regenerate exact-arithmetic checks.
The proof-status ledger distinguishes tested programs, compiled target
definitions, host proofs and embedded proof constructions.

\section{What is resolved, what remains, and how to use the package}
\label{sec:status}

\subsection{The mathematical results and their exact reach}

\begin{longtable}{P{.24\linewidth}P{.39\linewidth}P{.28\linewidth}}
\toprule
Problem & Result in this report & Remaining condition\\
\midrule\endhead
Compatible exact approximations & A finite uniqueness theorem gives
coherence of an explicit PR producer. & Correctness and uniqueness must
be proved at each level.\\[4pt]
Implicit algebraic construction & Filtered correction solves a simultaneous
polynomial system in each nilpotent truncation with a fixed step count.
& A seed root and inverse Jacobian at the first level.\\[4pt]
Formal power-series solutions & Coherent exact jets, with constant seed and
uniqueness. & Analytic radii and tail bounds are separate estimates.\\[4pt]
Local modular roots & Coherent solutions modulo \(q^n\), including composite
\(q\), from a unit Jacobian seed. & No inference to global integer or
rational solvability.\\[4pt]
Differential equations over integral rings & Hurwitz coefficients are
constructed by a triangular primitive recurrence without factorial
division. & Analytic interpretation needs additional assumptions.\\[4pt]
Geometric zero-count input & Sharp norm and subspace fiber bounds for
\(Y^m-cX^d\) with lower-weight terms and coprime \(m,d\). & The geometric
curve class is explicit; a general divisor library is not supplied.\\[4pt]
Finite spaces of functions & An explicit monomial basis and exact eventual
dimension \(L+1-(m-1)(d-1)/2\). & No identification with arbitrary
Riemann--Roch spaces is asserted.\\[4pt]
Finite-cover argument & A cover by \(M\) proper subspaces yields a bound
\(ML\) on a supplied distinct point list. & A proof or construction of the
cover is still needed.\\[4pt]
Deep arithmetic formalization & Compiled core and object composition;
explicit new targets and dependency order. & Full embedded lifting and
norm derivations remain unfinished.\\
\bottomrule
\end{longtable}

The norm certificate is produced from the coefficient tables. The jet and
modular root producers are constructed from their finite seeds. Thus the
applications do more than name a proof field containing the conclusion.
On the other hand, the finite-cover corollary is deliberately conditional:
its covering hypothesis is substantial mathematics, and a finite checker
for a proposed cover does not prove that an adequate cover always exists.

\subsection{Three boundaries that cannot be erased by a representation}

\paragraph{A count bound is not a cutoff.}
For a machine code \(e\), let \(R(e,t)\) say that the machine first halts
at exactly step \(t\). This is a PR predicate, and determinism proves
that at most one \(t\) satisfies it. If there were a total computable
bound \(B(e)\) containing every such witness, checking the first
\(B(e)\) steps would decide halting. Hence no uniform conversion from
an arbitrary finite-cardinality bound to a search cutoff exists, even
for PR families of sets with at most one element. This is a general
obstruction to an inference, not an independence result for any specific
Diophantine theorem. A new height theorem for a particular curve family
could still supply the missing information.

\paragraph{Finite Cauchy approximations do not decide exact algebraic cases.}
Define a real \(a_e\) to be \(2^{-t}\) if machine \(e\) first halts
at step \(t\), and zero if it never halts. At precision \(k\), simulate
only \(k\) steps and return the corresponding dyadic value if halting
has occurred, otherwise zero. This is a uniform PR rational evaluator;
its error is at most \(2^{-k}\). Externally, in the ordinary classical metatheory, every \(a_e\) is rational, yet a
uniform exact zero test for these names would decide halting. A claim
that every Cauchy-presented rational can be uniformly converted to an
exact rational code with decidable equality is therefore false.

Our regular lifting theorem uses a different input: an exact seed and
a finite invertibility certificate. For \(Y^2-X=0\) at \((0,0)\),
the derivative with respect to \(Y\) is zero, so the hypothesis fails.
The theorem does not silently select a singular Puiseux branch.

\paragraph{Coherent local roots need not be global points.}
The polynomial \(Z^2-2\) has the root \(3\) modulo \(7\), and its
derivative there is \(6\), a unit. The modular lifting theorem constructs
compatible roots modulo every \(7^n\). It has no rational root: writing
a putative root in coprime integers leads to both numerator and denominator
being even. Thus the tower supplies local arithmetic information, not an
integer or rational witness. Likewise a formal identity is not already
a quantified analytic error bound, even when an independent convergence
theorem is available for the same data.

\subsection{Consequences for the unresolved repository requests}

The new norm theorem settles the zero-count component of the Siegel
finite-cover route for the stated weighted curves. The dimension and
function-field constructions reduce further algebraic setup in that class.
The quantitative approximation theorem, suitable covers, local estimates,
and the global height/list problem remain. This does not prove the full
Siegel or Faltings theorem in \ipra.

The lifting theorem provides a common finite backend for the regular
parts of formal implicit construction, modular arithmetic and local
expansions. It does not construct all singular Newton--Puiseux branches
from arbitrary Cauchy coefficients. The existing exact-algebraic Puiseux
argument and analytic regular-chart estimates remain relevant separate
modules.

For Mordell--Weil the obstacle identified in the inspected packet is not
ordinary height descent. Its supplied representative list must still be
constructed and proved to cover the desired quotient. The checked descent
core in this archive can consume such data, but is not a resolution of
that missing arithmetic existence theorem. For Chebotarev,
Siegel--Walfisz and Bombieri--Vinogradov the analytic estimates and their
uniform witnesses similarly remain outside the results proved here.

\subsection{Reproducibility and the next concrete formalization steps}

The archive includes the complete LaTeX source and compiled PDF,
standalone Lean 4.32.0 sources, exact-arithmetic example programs, source
provenance and machine-readable validation records. The programs check
finite instances and invariants. Their passing results are regression
evidence, not proofs of the universal algebraic claims.

The first useful extension is to finish the common integer, rational,
list, polynomial and finite determinant library inside the object theory.
The next is the generic bounded-loop compiler and its coefficient-prefix
invariants. With those in place, frozen-inverse lifting has a short proof
dependency graph: polynomial Taylor remainder, Jacobian congruence,
residual improvement, uniqueness and restriction. The weighted norm
development then reuses the same polynomial kernel and adds a finite
determinant library, rescaling, Vandermonde factors and finite-list
grouping. Each endpoint is a \texttt{Prf PRAAx} object, transported to the
native calculus only afterwards.

The point of the framework is that these are common reusable modules.
The full logic of finite observation, the data consumed by a construction,
and the exact theorem remaining to be proved are visible at each boundary.
That makes future successes composable and prevents a failed existence
argument from disappearing inside a convenient representation.

\begin{thebibliography}{99}
\addcontentsline{toc}{section}{References and pinned source locations}
\bibitem{repoIndex} Vladimir Reshetnikov, ProveIt repository.
\emph{ProveIt proof-packet index}.
\href{https://github.com/VladimirReshetnikov/ProveIt/blob/052fbaebfc955294b4c02f6532ab7e7a3f760f7a/Logic/PrimitiveRecursiveArithmetic/docs/proofs/README.md}{Pinned source at commit \texttt{052fbaebfc95}}.
Exact path: \path{Logic/PrimitiveRecursiveArithmetic/docs/proofs/README.md}.

\bibitem{repoCore} Vladimir Reshetnikov, ProveIt repository.
\emph{ConciseIPRA: ordinary iPRA syntax and axioms}.
\href{https://github.com/VladimirReshetnikov/ProveIt/blob/052fbaebfc955294b4c02f6532ab7e7a3f760f7a/Logic/PrimitiveRecursiveArithmetic/ConciseIPRA/Lean/ConciseIPRA/PR.lean}{Pinned source at commit \texttt{052fbaebfc95}}.
Exact path: \path{Logic/PrimitiveRecursiveArithmetic/ConciseIPRA/Lean/ConciseIPRA/PR.lean}.

\bibitem{repoBridge} Vladimir Reshetnikov, ProveIt repository.
\emph{ConciseIPRA: native bi-interpretation}.
\href{https://github.com/VladimirReshetnikov/ProveIt/blob/052fbaebfc955294b4c02f6532ab7e7a3f760f7a/Logic/PrimitiveRecursiveArithmetic/ConciseIPRA/Lean/ConciseIPRA/NativeBiInterpretation.lean}{Pinned source at commit \texttt{052fbaebfc95}}.
Exact path: \path{Logic/PrimitiveRecursiveArithmetic/ConciseIPRA/Lean/ConciseIPRA/NativeBiInterpretation.lean}.

\bibitem{repoSmallFramework} Vladimir Reshetnikov, ProveIt repository.
\emph{Earlier compact framework packet}.
\href{https://github.com/VladimirReshetnikov/ProveIt/blob/052fbaebfc955294b4c02f6532ab7e7a3f760f7a/Logic/PrimitiveRecursiveArithmetic/docs/proofs/foundations/ipra_framework_2026-11-15/framework.tex}{Pinned source at commit \texttt{052fbaebfc95}}.
Exact path: \path{Logic/PrimitiveRecursiveArithmetic/docs/proofs/foundations/ipra_framework_2026-11-15/framework.tex}.

\bibitem{repoEmbedding} Vladimir Reshetnikov, ProveIt repository.
\emph{Finite-demand embedding framework; earlier hyperelliptic norm repair}.
\href{https://github.com/VladimirReshetnikov/ProveIt/blob/052fbaebfc955294b4c02f6532ab7e7a3f760f7a/Logic/PrimitiveRecursiveArithmetic/docs/proofs/foundations/ipra_embedding_framework/paper/05_diophantine.tex}{Pinned source at commit \texttt{052fbaebfc95}}.
Exact path: \path{Logic/PrimitiveRecursiveArithmetic/docs/proofs/foundations/ipra_embedding_framework/paper/05_diophantine.tex}.

\bibitem{repoFramework} Vladimir Reshetnikov, ProveIt repository.
\emph{Finite-contracts framework: scope and existing results}.
\href{https://github.com/VladimirReshetnikov/ProveIt/blob/052fbaebfc955294b4c02f6532ab7e7a3f760f7a/Logic/PrimitiveRecursiveArithmetic/docs/proofs/foundations/ipra_framework/paper/sections/01_scope.tex}{Pinned source at commit \texttt{052fbaebfc95}}.
Exact path: \path{Logic/PrimitiveRecursiveArithmetic/docs/proofs/foundations/ipra_framework/paper/sections/01_scope.tex}.

\bibitem{repoSiegel} Vladimir Reshetnikov, ProveIt repository.
\emph{Siegel packet: the unresolved general route and geometric fiber contract}.
\href{https://github.com/VladimirReshetnikov/ProveIt/blob/052fbaebfc955294b4c02f6532ab7e7a3f760f7a/Logic/PrimitiveRecursiveArithmetic/docs/proofs/diophantine-number-theory/siegel_ipra/paper/05_general_route.tex}{Pinned source at commit \texttt{052fbaebfc95}}.
Exact path: \path{Logic/PrimitiveRecursiveArithmetic/docs/proofs/diophantine-number-theory/siegel_ipra/paper/05_general_route.tex}.

\bibitem{repoMW} Vladimir Reshetnikov, ProveIt repository.
\emph{Mordell--Weil packet: actual quotient representatives and elliptic descent}.
\href{https://github.com/VladimirReshetnikov/ProveIt/blob/052fbaebfc955294b4c02f6532ab7e7a3f760f7a/Logic/PrimitiveRecursiveArithmetic/docs/proofs/diophantine-number-theory/mordell_weil_ipra/paper/03_elliptic.tex}{Pinned source at commit \texttt{052fbaebfc95}}.
Exact path: \path{Logic/PrimitiveRecursiveArithmetic/docs/proofs/diophantine-number-theory/mordell_weil_ipra/paper/03_elliptic.tex}.

\bibitem{repoFaltings} Vladimir Reshetnikov, ProveIt repository.
\emph{Faltings packet: cardinality, complete-list and height distinctions}.
\href{https://github.com/VladimirReshetnikov/ProveIt/blob/052fbaebfc955294b4c02f6532ab7e7a3f760f7a/Logic/PrimitiveRecursiveArithmetic/docs/proofs/diophantine-number-theory/faltings_ipra/faltings_ipra.tex}{Pinned source at commit \texttt{052fbaebfc95}}.
Exact path: \path{Logic/PrimitiveRecursiveArithmetic/docs/proofs/diophantine-number-theory/faltings_ipra/faltings_ipra.tex}.

\bibitem{repoGeometry} Vladimir Reshetnikov, ProveIt repository.
\emph{Diophantine geometry: modular obstructions and denominator data}.
\href{https://github.com/VladimirReshetnikov/ProveIt/blob/052fbaebfc955294b4c02f6532ab7e7a3f760f7a/Logic/PrimitiveRecursiveArithmetic/docs/proofs/diophantine-number-theory/diophantine_geometry_ipra/sections/06_obstructions.tex}{Pinned source at commit \texttt{052fbaebfc95}}.
Exact path: \path{Logic/PrimitiveRecursiveArithmetic/docs/proofs/diophantine-number-theory/diophantine_geometry_ipra/sections/06_obstructions.tex}.

\bibitem{repoPuiseux} Vladimir Reshetnikov, ProveIt repository.
\emph{Puiseux packet: exact algebraic input and Cauchy-name obstruction}.
\href{https://github.com/VladimirReshetnikov/ProveIt/blob/052fbaebfc955294b4c02f6532ab7e7a3f760f7a/Logic/PrimitiveRecursiveArithmetic/docs/proofs/algebra/puiseux_ipra/README.md}{Pinned source at commit \texttt{052fbaebfc95}}.
Exact path: \path{Logic/PrimitiveRecursiveArithmetic/docs/proofs/algebra/puiseux_ipra/README.md}.

\bibitem{repoLagrange} Vladimir Reshetnikov, ProveIt repository.
\emph{Lagrange inversion packet: finite truncation architecture}.
\href{https://github.com/VladimirReshetnikov/ProveIt/blob/052fbaebfc955294b4c02f6532ab7e7a3f760f7a/Logic/PrimitiveRecursiveArithmetic/docs/proofs/algebra/lagrange_inversion_ipra/README.md}{Pinned source at commit \texttt{052fbaebfc95}}.
Exact path: \path{Logic/PrimitiveRecursiveArithmetic/docs/proofs/algebra/lagrange_inversion_ipra/README.md}.

\bibitem{repoBell} Vladimir Reshetnikov, ProveIt repository.
\emph{Bell-number packet: finite identities and generating series}.
\href{https://github.com/VladimirReshetnikov/ProveIt/blob/052fbaebfc955294b4c02f6532ab7e7a3f760f7a/Logic/PrimitiveRecursiveArithmetic/docs/proofs/combinatorics/bell_numbers_ipra/paper/03_series.tex}{Pinned source at commit \texttt{052fbaebfc95}}.
Exact path: \path{Logic/PrimitiveRecursiveArithmetic/docs/proofs/combinatorics/bell_numbers_ipra/paper/03_series.tex}.

\bibitem{repoPrimes} Vladimir Reshetnikov, ProveIt repository.
\emph{Prime-distribution packet: exact scope and missing estimates}.
\href{https://github.com/VladimirReshetnikov/ProveIt/blob/052fbaebfc955294b4c02f6532ab7e7a3f760f7a/Logic/PrimitiveRecursiveArithmetic/docs/proofs/analytic-number-theory/prime_distribution_ipra/paper/sections/01_scope.tex}{Pinned source at commit \texttt{052fbaebfc95}}.
Exact path: \path{Logic/PrimitiveRecursiveArithmetic/docs/proofs/analytic-number-theory/prime_distribution_ipra/paper/sections/01_scope.tex}.

\bibitem{Hu} Yining Hu.
\emph{Explicit Implicit Function Theorem for All Fields}.
arXiv:1611.01415, 2016. \url{https://arxiv.org/abs/1611.01415}.
Cited for mathematical context; the finite proofs in this report are given directly.

\bibitem{AlonsoLombardiNeuwirth} Mar\'ia Emilia Alonso Garc\'ia,
Henri Lombardi, and Stefan Neuwirth.
\emph{Multivariate Hensel Lemma for ultrametric fields}.
Journal of Algebra 667 (2025), 803--823;
preprint \url{https://arxiv.org/abs/2301.06546}.
Cited for the constructive Hensel context, not as an unverified iPRA dependency.

\bibitem{StacksNorm} The Stacks Project Authors.
\emph{Trace and norm}, Section 9.20, tag 0BIE.
\url{https://stacks.math.columbia.edu/tag/0BIE}.
Cited for the determinant description of norms; the required finite identities
and weighted application are developed in the text.
\end{thebibliography}

\end{document}
